\documentclass[submitted,12pt]{article}
\usepackage{chemfig}
\usepackage{els}
\usepackage{cassirer}

\usepackage[version=4]{mhchem}

\newcommand{\IRt}{\bt{Identité et réalité}\xspace}
\newcommand{\IR}[2]{\citep[#1/#2]{Meyerson1908}}
\renewcommand{\SF}[2]{\citep[#1/#2]{Cassirer1910}}

\title{Cassirer's Philosophy of Chemistry}
\begin{document}

\maketitle

\begin{abstract}
\lipsum[1]
\end{abstract}

\section*{Introduction}

On \datef{30}{9}{1909}, during a visit to Paris, Hermann Cohen, the founder of the so-called Marburg school of neo-Kantianism, sent to Émile Meyerson a postcard asking whether he would be available to meet the following afternoon. A trained chemist, \citet{Meyerson1908} had just published his first philosophical book, \IRt , which brought him almost overnight prominence in Parisian intellectual circles at nearly fifty years of age \citep{Bensaude-Vincent2016a}. In the same postcard, Cohen noted that one of his former students, Ernst Cassirer, was also in Paris and would join them for the meeting. The following day, \datef{1}{10}{1909}, at around 2 p.m., they met at the Café de la Régence. A few weeks after the meeting, Meyerson received a copy of \EPt, \citets{Cassirer1906,Cassirer1907b} \letterecndp{Meyerson}{Cassirer}{2}{11}{1909}[121]\footnote{The numbering of the Meyerson--Cassirer letters follows \cite[\V{18}]{ECN}. However, the letters are not transcribed. Scans of the originals are available online at
\href{https://digitalisate.sub.uni-hamburg.de/recherche?tx_dlf_listview\%5Baction\%5D=main&tx_dlf_listview\%5Bcontroller\%5D=ListView&tx_dlf_listview\%5Bpage\%5D=1&tx_dlf_listview\%5Bsearch\%5D\%5Border\%5D=&tx_dlf_listview\%5Bsearch\%5D\%5BorderBy\%5D=&tx_dlf_listview\%5Bsearch\%5D\%5Bquery\%5D=Meyerson&cHash=0d009ea45cd8a5124ca05a24b99420c9}{Hamburger Kulturgut Digital}; last consulted on \today}. Interwoven with the warp of historical detail, the book traces the weft of a quasi-teleological development of modern science from \SB to \FB, an opposition that Cassirer had already begun to articulate systematically in his \emph{Habilitationsvortrag} \citep{Cassirer1907}.

By late December 1909, Meyerson had already completed the \s{first reading} of the \emph{Erkenntnisproblem}, and he had contacted Xavier Léon, the editor of the celebrated \emph{Revue de métaphysique et de morale}, with the intention of drafting a \emph{revue critique} Cassirer’s work \letterecndp{Meyerson}{Cassirer}{13}{12}{1909}[132]. From the outset, Cassirer welcomed Meyerson’s proposal with great enthusiasm \letterecndp{Meyerson}{Cassirer}{23}{12}{1909}[132]. Cassirer did not want to pressure Meyerson but he was clearly eager to hear his reaction before and joe that review was, on the second edition of the \EPt \hide{PBerlin 21. XII. 190 Berlin 17. I. 1910)}. Meyerson wrote to Cassirer a few months later, in April 19010, after completing the \s{second reading} of the \EPt, announcing a visit to Berlin. Between \datef{25}{4}{1909} and \datef{3}{5}{1909}, Cassirer welcomed Meyerson into his home on more than one occasion. The correspondence does not reveal the content of their conversations. Nevertheless, it seems legitimate to suppose that Meyerson anticipated the criticism that would later appear in Meyerson's published review. In particular, Meyerson might already have objected that \emph{Das Erkenntnisproblem} had completely neglected the role of chemistry in the historical development of the \s{problem of knowledge} in early modern science. It was Cassirer’s exclusive focus on the history of  physics that had led him to overemphasize the role of function-concepts at the expense of the substance-concept. 

%Natorp evidently did not consider chemistry worthy of discussion in a book devoted Indeed, chemistry is conspicuously absent from Natorp's \citetitle{Natorp1910a}, the other major work of the Marburg school published that same year \citeyear{Natorp1910a}. Natorp evidently did not consider chemistry worthy of discussion in a book devoted to the foundations of the \s{exact sciences}. to the foundations of the \s{exact sciences}. 

%For Cassirer precisely the opposite is the case. Chemistry does share the same model of concept formation \origg{Begriffsbildung} as physics, but at the same time, it sheds new light on it. Unlike physics, which aims at general laws at the nature of individual substances as in formulas. \todo{improve}.

At that time, Cassirer was already in the final stages of drafting \SFt \citep{Cassirer1910}, whose completion followed not much later, around July 1910. Cassirer immediately asked his publisher to send at least part of the galley proofs to Meyerson \letterecndp{Meyerson}{Cassirer}{17}{7}{1910}[148], who, however, received the entirety of the book much later, when the book was already in press \todo{letter}. Meyerson was aware, that without having access to it he might have misunderstood Cassirer's stance in the \EPt. Indeed, he would have definitively realized that, far from neglecting chemistry, Cassirer had chosen to dedicate a rather the 20-page section \rom{8} of the fourth chapter to the topic of concept formation in the natural sciences \SF{**}{**}\footnote{The number following the slash indicates the corresponding page in the English translation given alongside the original edition in the bibliography}. Cassirer’s interest for chemistry is far from obvious. Chemistry is conspicuously absent from\citetitle{Natorp1910a}, the major work of the Marburg school that Paul Natorp concluded at about the same time in March \citeyear{Natorp1910a}. For Cassirer, by contrast, chemistry has a peculiar epistemological status with respect to physics among the exact sciences. In both cases, scientific formation proceeds from \SCs to \FCs; however, while physics foregrounds the movement toward increasingly general laws, only chemistry brings into view the complementary movement toward increasingly determinate individuality \citep[270–272; 289–293]{Cassirer1910}. 

Meyerson's review finally appeared in the January issue of the 1911 volume of \emph{Revue}. Cassirer read it immediately and conceded that from a \emph{historiographical} perspective, \EPt had given less space to the history of chemistry than was necessary \todo{letter}; indeed, he immediately added an entire new section dedicated to chemistry addressing Meyerson's complaint to the second edition of the \EPt which appear shortly thereafter \citep{Cassirer1911a}. However, from a \emph{systematic} point of view, Cassirer integrated this material into a very different conceptual framework than Meyerson had in mind. What was at stake was ultimately not the relevance of the history of chemistry, but the philosophical conclusions to be drawn from it. On Meyerson’s account, the history of chemistry, in contrast to that of physics, unmistakably shows that scientific explanation proceeds by progressively eliminating distinctions until matter is reduced to an undifferentiated unity \todo{improve}. For Cassirer, by contrast, the history of chemistry shows more clearly than that of physics that \s{generality} required of scientific concept-formation is not an end in itself but serves merely as a means of attaining increasing individual determinacy. Indeed the entire first part of \SFt Cassirer pointed out concluded the It is chemistry that challenges Heinrich \citets{Rickert1902} claim that individuality represent the \s{limit of scientific concept-formation}. 

Reflections on chemistry, from Duhem to Meyerson, played a central role in early French philosophy of science \citep{Bensaude-Vincent2009a}. By contrast, German philosophy, perhaps following Kant's own denial of chemistry's status as a \s{science}, had always remained physics-oriented. Cassirer's interest of chemistry represents a unique exception. However, despite the increasing attention given to his philosophy of mathematics and physics, it has received no attention. This paper aims to fill a gap in the literature by reconstructing that background. While Meyerson left an important legacy in France to Metzger and Bachelard, Cassirer’s reflections on chemistry seem to have remained isolated in the German context, where even logical empiricism paid no attention to chemistry\footnote{Interestingly, \citet[108]{Schlick1932a} dismissed Meyerson's philosophy as the philosophy of chemists: \q{the psychological explanation of [Meyerson's] view lies in the fact that he started as a chemist, who is used to thinking in terms of identical substances}}.



%While physics seeks to establish general laws across a manifold of individual cases, chemistry seeks to determine the identity of individual elements or compounds 
%
%It is only through chemistry that the question of the relationship between the particular and the universal in scientific concept formation comes to light. 

%If, in both cases, science increasingly proceeds from substance-concepts to function-concepts, 


%Physics has been defined as the study of the properties of matter that are shared by all substances, whereas chemistry has been defined as the study of the properties of individual substances.

%Physics foregrounds the relations that hold across a manifold of particular cases; chemistry foregrounds the relations that determine the identity of a particular substance. Laws and formulas both express relations, but the former primarily articulate general dependencies, whereas the latter articulate specific chemical identities.

%Your point can therefore be stated precisely: chemical formulas are the counterparts of physical laws insofar as they provide relational determination; their distinctive achievement is to articulate the structure and identity of substances, whereas physical laws articulate causal dependencies among processes. The possible passage from the former to the latter is a further scientific achievement, rather than something already accomplished by writing a formula.

%Formulas articulate the determinations through which compounds are distinguished. Compositional formulas specify which elements occur and in what proportions; structural formulas differentiate compounds that share the same composition.


%Physics foregrounds the movement toward increasingly general laws; only chemistry brings into view the complementary movement toward increasingly determinate individuality.

%Physics foregrounds the generality of laws; chemistry foregrounds the specific determination of substances. 

% and how its characteristic properties belong together.

%While physics searches for general laws, chemistry aims at the determination of individual substances.  %what distinguishes one substance from another,

%complètement abstraction de ses théories


%who, however, received the entirety of the book much later, when the book was already in press \todo{letter}.


%c'est la prédominance nécessaire du concept de règle ou de fonction mathématique sur celui d'essence ou de substance, allant jusqu'à l'annihilation définitive du second au profit du premier

%peu près 
%chimiques et de ce qu'on pourrait en tirer pour sa conception de la
%matière et se borne à étudier ses opinions philosophiqu

\todo{abstract and generic concept of matter}
\todo{ way of dealing with individual substances that transcends
their singularities in making them samples of pure chemical species.}




%[Needham, 2004] P. Needham. Has Daltonian Atomism Provided Chemistry with any Explanations? Philosophy of Science, 71, 1038-47, 2004.
%[Needham, 2008] P. Needham. Resisting Chemical Atomism: Duhem’s Argument, Philosophy
%of Science, forthcoming.
%[Needham, 1981] P. Needham. Temporal Intervals and Temporal Order, Logique et Analyse, 93,
%49-64, 1981.
%[Needham, 2000] P. Needham. What is Water? Analysis, 60, 13–21, 2000.
%[Needham, 2002] P. Needham. The Discovery that Water is H2O’, International Studies in the
%Philosophy of Science, 16, 201–222, 2002.
%[Needham, 2006a] P. Needham. Aristotle’s Theory of Chemical Reaction and Chemical Sub-
%stances, in Davis Baird, Eric Scerri and Lee McIntyre, eds., Philosophy of Chemistry: Syn-
%thesis of a New Discipline, Boston Studies in the Philosophy of Science, Springer, Dordrecht,
%2006, 43-67, 2006.
%[Needham, 2006b] P. Needham. Substance and Modality, Philosophy of Science, 73, 829-40,
%2006.
%[Needham, 2007] P. Needham. Macroscopic Mixtures, Journal of Philosophy, 104, 26-52, 2007.
%[Needham, 2011] P. Needham. Microessentialism: What is the Argument? Nous, 45, 1-21,
%2011.
%[Needham, 2002] P. Needham. Duhem’s theory of mixture in the light of the Stoic challenge to the Aristotelian conception, Studies in the History and Philosophy of Science, 33: 685-708, 2002.


%For the primary task of  Insofar 

\todo{rare excepton}

%whereas the physicist, who goes
%more deeply into the explanation of nature, has to think in termsof events.

\todo{Entsoffliching Schummer}

\todo{Chemilca sace}




%never developed in a to a comaorable to french radio, inded, not but also loglca emproco no inrest in chemis.  Here is a major and special feature of modernphysical science: it is less a science of facts than a science of effects.


\section{Chemistry's Concept-Formation}
%

%Das Problem der Bej^iffsbildung in der Chemie.—Die Chemie
%der sinnlichen Qualitativ».; diePhlogistontheorie.—DasGesetzder
%bestimmten Proportionen >(J. D. Richter).—Daltons Gesetz der
%multiplen Proportionen. -^ Die Entwicklung des chemischen
%Atombegriffs.—Der Atomdjj^iff als Verhältnisbegriff.—Der
%„regulative" Gebrauch des A;.'vtnbegriffs

%According to Cassirer, the account of concept formation in the exact natural sciences remains incomplete from a logical standpoint as long as it fails to consider the fundamental concepts of chemistry alongside those of physics.

Cassirer referred to chemistry throughout \SF, but devoted section \rom{8} of the fourth chapter specifically to chemical concept formation \origg{Begriffsbildung)}. Chemistry had initially developed as a purely \s{descriptive science}, much like zoology or botany. However, with the emergence of physical chemistry in the late nineteenth century, it had firmly established itself among the \s{constructive sciences}\footnote{In support of this point, he could invoke Walther \citet{Nernst1896}, one of the founding fathers of modern physical chemistry, who had explicitly articulated this view}. In particular chemists increasingly turned to thermodynamics to investigate chemical processes, analyzing reaction rates and equilibrium conditions in terms of macroscopic variables such as temperature, pressure, and energy. As Cassirer points out, from this point of view, chemistry's \qt{fundamental laws, as for example the law of phases (\emph{Phasenregel}) of Gibbs or the law of chemical masses (\emph{Massenwirkung}), belong to the same purely mathematical type as the propositions of theoretical physics}{Ihre Grundgesetze, wie etwa die Gibbs'sche Phasenregel oder das Gesetz der chemischen Massenwirkung, gehören demselben rein mathematischen Typus an, wie irgendwelche Sätze der theoretischen Physik} \SF{271}{204}\footnote{While the \s{phase rule} concerns the coexistence of phases and the \s{law of mass action} specifies the equilibrium relations among reacting species, both receive a unified formulation within the thermodynamics of Gibbs, Helmholtz, Duhem, and Planck through the notion of thermodynamic potential. Alongside this program of determining the conditions of chemical equilibrium, molecular and ionic accounts of chemical phenomena were already being developed, notably by van 't Hoff and Arrhenius, without yet relying on a theory of the internal electronic structure of atoms \citep{Barkan1999}}. The difference between physics and chemistry thus lies not in their status as \s{exact sciences}, but in the historical paths by which they attained that status. Whereas physics had been a \s{constructive} discipline from the outset, chemistry acquired this character only at a later stage in its development. 


 %Physics has been defined as the study of the properties of matter that are shared by all substances, whereas chemistry has been defined as the study of the properties of individual substances.

%Physics foregrounds the relations that hold across a manifold of particular cases; chemistry foregrounds the relations that determine the identity of a particular substance. Laws and formulas both express relations, but the former primarily articulate general dependencies, whereas the latter articulate specific chemical identities.

%Your point can therefore be stated precisely: chemical formulas are the counterparts of physical laws insofar as they provide relational determination; their distinctive achievement is to articulate the structure and identity of substances, whereas physical laws articulate causal dependencies among processes. The possible passage from the former to the latter is a further scientific achievement, rather than something already accomplished by writing a formula.

%Formulas articulate the determinations through which compounds are distinguished. Compositional formulas specify which elements occur and in what proportions; structural formulas differentiate compounds that share the same composition.


%Physics foregrounds the movement toward increasingly general laws; only chemistry brings into view the complementary movement toward increasingly determinate individuality.

%Physics foregrounds the generality of laws; chemistry foregrounds the specific determination of substances. 

% and how its characteristic properties belong together.

%While physics searches for general laws, chemistry aims at the determination of individual substances.  %what distinguishes one substance from another,

%complètement abstraction de ses théories
%Der Werdegang einer Wissenschaft


%We might say that classical mechanics is primarily organized around states, quantities, and their evolution, whereas much of traditional chemistry is organized around substances and their possible transformations.

%The periodic law expresses regularities among the properties of different elements. It does not describe a chemical transformation, nor does it specify the behavior of a particular physical system.

%The periodic law does not merely summarize similarities already observed; it also allows scientists to infer properties of substances not yet investigated.

%The history of chemistry shows that scientific systematicity need not consist exclusively in the deduction of particular phenomena from increasingly general laws. It may also emerge through the progressive coordination of different empirical domains into a common relational system.

%Classical mathematical physics has privileged laws relating measurable quantities across broad classes of systems. Chemistry has developed quantitative laws of composition, regularities governing transformations of particular substances, and principles establishing systematic relationships among different substances.



Physics had the advantage of dealing with relatively simple systems, such as planetary motion, in which bodies could be treated as \s{structureless} mass points. It could therefore replace the thing-concept \origg{Dingbegriff} with the concept of law \origg{Gesetzesbegriff} at an early stage. Chemistry, by contrast, had to begin with the problem of the individual thing \origg{Einzelding} in all its complexity \SF{271}{204}---specific properties if properties of elements and compounds. Cassirer therefore asks whether chemistry requires a different form of concept formation from that of physics. As one might expect, Cassirer argues that this question can be answered only \qt{if we follow the concrete historical development of chemical doctrines themselves}{wenn man dem konkreten geschichtlichen Gang der chemischen Lehren selbst folgt} \SF{272}{204}. Given the systematic aims of the book, Cassirer does not attempt to trace the history of chemistry in detail, but seeks instead to identify \qt{the great logical lines of direction}{die großen logischen Richtlinien} along which it developed \SF{272}{204}. \todo{sources} Cassirer does not appear to draw directly on the original sources, relying instead on \citet{Wurtz1879}, \citet{Meyer1884}, \citet{Ladenburg1902}, \citet{Duhem1902} and \citet{Ostwald1906Leitlinien} as his principal historical accounts of the development of chemistry. He also draws on \citet{Berthelot1885} in his discussion of the transition from alchemy to qualitative chemistry. Cassirer reorganizes this material to fit his own philosophical narrative. Like the history of physics, the history of chemistry can be understood as a process of progressive \s{arithmetization}, in which qualitative differences are expressed in terms of numerical values and the functional relations between them. In chemistry, however, this process serves a different purpose: not the discovery of general laws, but the identification of individual substances.





%Laws of behavior establish how a system changes as a function of its state and the conditions under which it operates.
%Laws of order establish systematic relations among different objects, substances, or properties.


%Physics primarily investigates how systems behave, whereas chemistry primarily investigates what kinds of substances exist, how they are related, and how they can be transformed into one another.



%As one might expect, he places particular emphasis on the precise nature of the transition from the genus-concept (\emph{Gattungsbegriff}) to the sequence-concept (\emph{Reihenbegriff}). Nevertheless, the case of chemistry also brings a distinctive dimension to this narrative, particularly by clarifying the relationship between the universal and the individual.




%we dealing here merely with a gap that can be filled by newly addedeterminations, which nonetheless belong to the same logical direction of thought, or must a fundamentally different basic form of knowledge be recognized and introduced at this point? 

%offers no handle for this new problem. For here it is only the symbol of a particular form of connection, one that has increasingly stripped away every particular material content; it designates only the manner of a possible coordination, not the | "what" of the elements that are to be coordinated with one another. 



%Initially, \qt{the \s{concept}, in the specific meaning it has in mathematics and physics, is not available for this new problem}{Der „Begriff" aber enthält gerade in der spezifischen Bedeutung, die er in Mathematik und Physik besitzt, keine Handhabe für dieses neue Problem} \SF{271}{204}.  

%Yet the "concept," precisely in the specific sense it possesses in mathematics and physics, offers no means of addressing this new problem. For here it is merely the symbol of a particular form of connection, from which all specific material content has increasingly been stripped away; it designates only the manner of a possible correlation, not the




\subsection{From Richter to Dalton}

Initially, concept-formation in chemistry is dominated by the classical procces of abstraction of Aristotelian logic: \qt{The properties common to a class of things, which thus provide the basis for forming a particular \ls{genus-concept}, are at the same time singled out as physical constituents and endowed with independent existence}{Die Eigenschaften, die einer Klasse von Dingen gemeinsam sind und die somit die Grundlage für die Bildung eines bestimmten Gattungsbegriffs abgeben, werden zugleich als physische Bestandteile abgesondert und zu selbständiger Existenz erhoben} \SF{**}{**}. \cop{Solid bodies differ from liquids and volatile substances by virtue of a specific, separable quality: a change in their state of aggregation entails the disappearance of one quality and the emergence of an entirely different material nature. Mercury, for instance, could supposedly be converted into gold by removing the two \s{elements} responsible for its liquidity and volatility and substituting other qualities in their place.} In traditional chemistry, elements were conceived not merely as possessing characteristic \s{qualities}, but as their very embodiments. Sulfur represented combustibility, salt solubility, and mercury the characteristic properties of metals. Chemical reactions governed by regularities were thus explained by attributing the corresponding qualities to substantial entities. According to Cassirer, this conception of nature persisted into the modern period, most notably in Bacon's physics\todo{Back in EP}. 

In Cassirer’s reconstruction, concept formation in early chemistry followed the traditional model of abstraction, moving from the genus-concept to an ontological conception of substance. Chemists proceeded by reducing complex objects of sensory experience to a set of simple, abstract qualities in terms of which their properties could be explained \SF{**}{**}. Heat, for example, was regarded as a distinct entity whose presence produced the characteristic effects of warm bodies. Similarly, with the introduction of phlogiston, combustibility, a property that we seem to perceive directly in various bodies, was conceived as a distinct substance admixed with those bodies: \qt{[A]nd from this assumption, the whole structure of chemistry before Lavoisier follows with inner necessity}{Und diese Annahme ist es, aus welcher der gesamte Aufbau der Chemie vor Lavoisier mit innerer Notwendigkeit folgt} \SF{205}{155}. In a footnote to this passage, Cassirer refers to the \s{excellent presentation} in \citets{Meyerson1908} book, which assigns Lavoisier a pivotal role in the quantitative turn in the history of chemistry. Cassirer may well have added this reference at a late stage of revision. Indeed, his account immediately diverges from Meyerson's, assigning a central role instead to the German chemist Jeremias Benjamin Richter. Meyerson, by contrast, does not even mention Richter, although the latter figures prominently in Cassirer's principal historical sources \citepp[chap. III]{Duhem1902}[lect.~4]{Ladenburg1902}[I.II]{Wurtz1879}[**]{Ostwald1906Leitlinien}.
 
%They have the object of representing to us laws in accordance with which certain phenomena are produced, whereas chemical theory seeks to classify substances. 
 
In Cassirer's view, it was not Lavoisier's law of \s{conservation of mass}, but Richter's \s{law of definite or constant proportions} that constituted \qt{the starting-point of modern chemical \ls{theory}}{ersten Ausgangspunkt der modernen chemischen Theorie} \SF{273}{205} and marked a decisive break with the alchemical tradition. Richter's investigations into chemical neutralization helped lay the foundations of quantitative chemistry by establishing, at least implicitly, two principles: that substances combine in fixed quantitative ratios (definite proportions), rather than in the variable proportions characteristic of physical mixtures, and that these ratios are systematically related across different chemical combinations (reciprocal proportions)\footnote{The attribution varies across historical accounts. Ostwald, for example, attributes the law of definite proportions to Proust and the law of reciprocal proportions to Richter}\todo{check}. According to Cassirer, as in physics, \qt{[t]he demand for strict numerical determination here also forms the decisive turning point}{Die Forderung strenger Zahlbestimmungen bildet auch hier den entscheidenden Wendepunkt} \SF{272}{205}. What matters for Cassirer is not that the permanence of substance takes the form of the conservation \origg{Beharrung} of a quantity as in Lavoisier, but rather the establishment of a coordination \origg{Zuordnung} between two distinct ordered sequence of numbers. 

At this stage, one can establish these laws by measuring the masses of substances involved in chemical reactions \qt{entirely independently of any conception of the constitution of matter, and in particular independently of the atomic hypothesis}{völlig unabhängig von jeder Auffassung über die Konstitution der Materie, insbesondere unabhängig von der atomistischen Hypothese} \SF{273}{205}. Following \citet[71]{Duhem1902}, Cassirer considers a sequence of acids and a sequence of bases 

\begin{equation*}
A_1, A_2, A_3, \ldots \qquad B_1, B_2, B_3, \ldots\,. 
\end{equation*}
%
Each of these substances can be put into correspondence with one of the numbers in the sequence

\begin{equation*}
m_1, m_2, m_3, \ldots \qquad n_1, n_2, n_3, \ldots\,,
\end{equation*}
%
so that every acid $A_i$ is assigned a characteristic number $ m_i,$ and every base $B_j$ another characteristic number $ n_j.$ Suppose sulfuric acid is assigned the number \(m_1=100\). If 100 grams of sulfuric acid are neutralized by 82 grams of sodium hydroxide, sodium hydroxide is assigned \(n_1=82\). If 82 grams of sodium hydroxide are in turn neutralized by 74 grams of hydrochloric acid, hydrochloric acid is assigned \(m_2=74\). Each substance is thus characterized by a number expressing its quantitative combining relations. Once $m_p$ and $n_q$ are known, the relation for the pair $A_p, B_q$ follows: \qt{The manner, in which an element of the first sequence combines with an element of the second, is univocally \origins{eindeutig} determined by these numbers; the two weights, according to which an acid $A_p$ combines with any base $B_q$, are related according to the corresponding numerical values $m_p$ and $n_q$}{Die Art, in der sich irgend ein Element der ersten Reihe mit einem Element der zweiten verbindet, ist durch diese Zahlen eindeutig bestimmt: die beiderseitigen Massen, nach welchen irgend eine Säure $A_p$ mit irgend einer Basis $B_q$ zusammentritt, verhalten sich wie die entsprechenden Zahlwerte $m_p$ und $n_q$} \SF{273}{205}. 

Richter tabulated the quantities of the bases which are neutralized by the same weight of sulphuric acid, of hydrofluoric acid, etc., and these he calls neutrality sequence, or sequence of masses; he also determined the quantities of the acids which are saturated by. the same quantity of different bases. In doing this, he thought he had discovered certain regularities. The sequence of masses in the case of the bases formed an arithmetical, and that in the case of the acids, an arithmetic progression. This particular hypothesis turned out to be empirically false \citepp[**]{Ladenburg1902}[**]{Wurtz1879}[**]{Ostwald1906Leitlinien}. Like \citet[144]{Duhem1902}, however, Cassirer aims to emphasize that Richter borrowed nothing from atomistic doctrines; his numerical speculations instead expressed the same \qt{general Pythagorean doctrine of the \s{harmony} of the cosmos}{allgemeine Pythagoreische Grundansicht von der „Harmonie" des Alls} \SF{273}{206} found in Kepler's work. In Cassirer's characteristic rhetoric, this case illustrates how \cop{scientific concept-formation does not primarily consist in discovering what each of these things \s{really is} internally. It consists instead in finding a rule that assigns each member a determinate place within an ordered system}.\todo{improve}


%Die Fassung, die das Gesetz der konstanten Verbindungszahlen bei seinem eigentlichen  wissenschaftlichen Begründer erhält, fügt allerdings dieser allgemeinen Grundanschauung  sogleich einen neuen konkreten Zug hinzu. Sachlich ist auch hier zunächst nur ausgesprochen, daß es für jedes Element eine charakteristische Äquivalenzzahl gibt und daß, wenn je zwei oder mehrere Elemente in eine Verbindung eingehen, ihre Massen sich wie ganze Vielfache  dieser Grundzahlen verhalt


%Richter merely claimed that for every element there is a \qt{characteristic equivalence-number}{charakteristische Äquivalenzzahl} \SF{274}{206}, that is, an equivalent weight. 

%compounds with definite compositions from mixtures of variable composition

%One compound always has the same composition. Different compounds formed from the same elements have compositions related by simple integer ratios.

%Small whole-number ratios




%Richter's law makes chemical combining weights part of a coherent system of relations. Once a reference element is chosen, the combining weights of other elements can be compared systematically.

%but the possibility of organizing chemical combinations into a coherent system of quantitative relationships.

As the conventional rational recontruction goes, Richter's scheme ran into a difficulty: the same two elements can form more than one compound, each with its own definite composition. In other terms,  they do not couple in one unique set of proportions. Carbon and oxygen, for example, form both carbon monoxide and carbon dioxide. In the former, $12\text{g}$ of carbon combine with $16\text{g}$ of oxygen; in the latter, $12\text{g}$ of carbon combine with $32\text{g}$ of oxygen. A single proportional number for oxygen cannot account for this variation. It was John Dalton to provide a solution to the problem formulating the \s{law of multiple proportions}. Instead of discarding the idea of characteristic masses $\left(p_i\right)$, he retained them as invariant base units and introduced integer multipliers ($\lambda, \mu, \nu$). The mass ratio of any compound is then no longer simply $p_1: p_2$, but

\begin{equation*}
\lambda p_1 : \mu p_2 : \nu p_3 : \cdots
\end{equation*}
%
where $p_i$ is the invariant characteristic mass of the element (e.g., $p_C=12$, $p_O=16$) and $\lambda, \mu, \nu$ are whole numbers ($1, 2, 3, \ldots$). In carbon monoxide, the mass ratio is $1 \cdot 12 : 1 \cdot 16$, so that $\lambda=1$ and $\mu=1$; in carbon dioxide, it is $1 \cdot 12 : 2 \cdot 16$, so that $\lambda=1$ and $\mu=2$. This is the origin of the Berzelius's \emph{compositional} formulas, which express the proportions in which elements combine to form compounds. On this purely quantitative reading, a formula such as 

\begin{equation*}
\mathrm{H}_2\mathrm{O} \quad \text{or} \quad \mathrm{CO}_2
\end{equation*}
%
is an algebraic shorthand for these integer multipliers. The letters represent the characteristic proportional masses $\left(p_1, p_2\right)$, and the subscripts the integers ($\lambda, \mu$). \cop{Formulas of this kind allow not only to summarize given empirical experiences, to elicit new ones}: \cop{Arranging their formulas into a numerical sequence might expose a gap: oxygen occurs in quantities of 1, 3, and 4, suggesting a compound with 2}.\todo{Louis-Joseph Gay-Lussac}
\todo{rational and not irrational}

\todo{The compositional formula H2O is an ntirely macroscopic concept, reflecting the gravimetric proportions of hydrogen to oxygen in water of 1 : 8. It is scaled in such a way as to provide a sense to saying that cuprous oxide, for example, with compositional formula Cu2O, contains just as much copper as water contains hydrogen, despite its composition of 127 parts by weight of copper to 16 of oxygen; and hydrogen chloride, with a
compositional formula HCl and gravimetric proportions 1 : 35.5, contains half as much hydrogen as does water.}

%The assumption of definite proportions, which forms the basis for the stoichio- metric compositional formulas on which the 19th century conception of structuralformulas differentiating isomeric variations was built, didn’t survive unscathed into the 20th century.

%Cassirer's contrast may be between laws that determine possible processes and formulas that determine particular substances
\todo{laws vs. formulas}
\todo{Duhem main source that postiiv to exit, to deny that atoms are substace.}
\todo{no return to teh Aristotelian miste}

The goal \q{law of constant combining numbers} \origg{Verbindungszahlen} was initially to characteristic equivalent number for each element \origg{Äquivalenzzahl} and that whenever two or more elements enter into a compound, their masses stand to one another \qt{as whole multiples of these basic numbers}{wenn je zwei oder mehrere Elemente in eine Verbindung eingehen, ihre Massen sich wie ganze Vielfache  dieser Grundzahlen verhalt} \SF{**}{**}. In principle, one can establish the laws of definite, reciprocal and multiple by measuring the macroscopic masses of substances involved in chemical reactions. In Dalton's work, however, this rule of \s{multiple proportions} immediately merges with a particular interpretation: \qt{The concept of combination-weight \origins{Verbindungsgewichts} is transformed into that of atomic weight \origins{Atomgewichts}}{Der Begriff des Verbindungsgewichts wandelt sich zu dem des Atomgewichts} \SF{274}{206}. If matter were continuous, like a fluid, elements should be able to combine in any proportion whatsoever. However, according to the law of multiple proportions, 12 g of carbon combine with exactly 16 g or exactly 32 g of oxygen. The multipliers ($\lambda$, $\mu$, $\nu$) are whole numbers; there is no intermediate possibility. Dalton reasoned that if substances always combine in fixed mass proportions, it must be because they consist of indivisible units, namely atoms, with precisely those relative atomic weights. Different elements possess atoms of different masses, while all atoms of the same element possess the same invariant mass: \qt{Statements concerning the essential nature of the ultimate constituent parts of individual bodies have replaced the empirically determined numerical proportions \origg{Proportionalzahlen} of those bodies}{An die Stelle  der empirisch ermittelten Proportionalzahlen der einzelnen Körper sind Aussagen über eine  wesentliche Beschaffenheit ihrer letzten konstitutiven Bestandteile getret} \SF{**}{**} Dalton suggest that chemical mass proportions could be used to calculate relative atomic weights by taking Hydrogen $= 1$ as the unit standard. 

%Dalton realized that these integer relationships could be explained if matter is composed of discrete, indivisible units (**atoms**), each characterized by a constant relative **atomic weight**


%imple substances are formed of atoms. The atoms of the same simple substance
%are all alike; they have the same mass. The atoms of two different simple substances
%have different masses, and their masses stand to one another as do the equivalents of
%simple substances. Thus, for the various simple substances, the equivalents measure
%the atomic masses and can be attributed this name.
%All compound substances are reducible to molecule

%The inference of the "chemical atom" arises specifically because the macroscopic laws yield **discrete, whole-number constants**

%**whole multiples** (*ganze Vielfache*) of a characteristic equivalence number, the concept of combining weight naturally transforms into the concept of a discrete **atomic weigh**
\todo{stoichiometric ratio }
\todo{Wallston might prove Cassirer point: using atoms forces A system founded strictly on equivalents had to invent different equivalents for the same element depending on its role, destroying systematic unity }

\todo{While an element may have several equivalent combining weights depending on the reaction, it possesses only one unique, invariant relative atomic weight}

%The case of Wollaston demonstrates that macroscopic laboratory ratios alone cannot produce a unified, unambiguous chemical system [47–49, 55]. Without the atomic hypothesis (and the physical laws that verify it), equivalents collapse into a conflicting collection of ratios. The atomic hypothesis was necessary because it provided an invariant standard that reconciled multiple proportions, gas densities, and structural isomerism into a coherent whole



%The concept of a single, immutable equivalent weight fails because elements combine in multiple proportions. For instance, Jean-Baptiste Dumas highlighted that copper possesses different equivalent (or combining) weights depending on whether it forms cuprous or cupric oxide.

Also, in the next stage, Cassirer follows the conventional account. Dalton’s hypothesis in turn faced a fundamental difficulty. The empirical facts do not determine the relative atomic weights uniquely. Suppose hydrogen is taken as the unit: the same experimental composition can be represented using different atomic weights for oxygen, provided that the molecular formulas are altered accordingly. Thus the experimental facts of composition alone cannot decide between, for example, $O=16$ and $O=8$. The same applies to sulfur, for which values such as $ S=8, \quad S=16, \quad S=32, \ldots$ are equally admissible \SF{274--75}{206--7}. Chemists could have accepted this ambiguity as a fact; instead, they introduced additional, independent relations to determine the ratios unambiguously. Following \citets{Meyer1884} account, Cassirer points out how the \emph{univocal} \origg{Eindeutig} the relative atomic weights is progressively determined not by a single measurement, but by the mutual convergence and correction of independent physical laws. Cassirer, following Meyer, mentions three especially important ones: Avogadro’s rule\footnote{The Principle: Equal volumes of all gases under the same pressure and temperature contain equal numbers of molecules ($M \propto d$). By measuring vapor densities across a wide array of volatile compounds containing a given element, chemists determine the smallest mass quantity of that element present in two volumes of gas. This sets the definitive upper boundary for the element's atomic weight }, connected with gas volumes and vapor densities; the Dulong--Petit law, connecting atomic weights with heat capacities\footnote{The product of an element's solid specific heat and its atomic weight is approximately constant ($A \cdot c \approx 6.0\text{–}6.4$): For non-volatile solid elements—especially metals—where vapor densities cannot be measured, Dulong and Petit's law provides a direct thermal determination of atomic weight}; and Mitscherlich’s law of isomorphism, connecting crystalline form with chemical composition and structure\footnote{Compounds possessing identical crystalline forms contain the same number of atoms combined in the same structural Application: Where neither vapor density nor specific heat is accessible, an unknown atomic weight can be deduced by substituting the element into an isomorphic matrix of a known compound (e.g., perchlorates and permanganates)} \SF{275}{207}. No single rule can univocally determine the relative atomic weights; their strength lies in their capacity to correct one another's anomalies. Relative atomic weights become increasingly determinate because several initially distinct sequence of phenomena are made mutually consistent: \qt{It is only the totality of all these different points of view, mutually confirming and correcting each other, that finally after many experiments gives a unitary table of atomic weights, and thus lays the basis of a definite system of chemical formulae}{Erst die Gesamtheit aller dieser verschiedenen Gesichtspunkte, die sich wechselseitig bestätigen oder korrigieren, ergibt schließlich nach mannigfachen Versuchen eine einheitliche Tabelle der Atomgewichte und legt damit den Grund zu einer eindeutigen chemischen Formelsprache} \SF{275}{207}.


For Cassirer it becomes clear how the \qt{the conditions of scientific \ls{production} differ from those of critical \ls{reflection}}{Die Bedingungen der wissenschaftlichen Produktion sind andere, als es die der kritischen Reflexion sind} \SF{**}{**}. A modern chemist might maintain that \qt{[t]he objective existence of the different types of atoms is presupposed; it is only necessary to discover their \emph{properties}, and to define them more exactly}{Die objektive Existenz der verschiedenen Arten von Atomen wird vorausgesetzt: nur ihre Eigenschaften gilt es noch zu ermitteln und quantitativ schärfer zu umgrenzen} \SF{275}{207}. As inquiry advances and encompasses an increasingly diverse range of phenomena, atomic properties appear to become richer and more determinate, as though the underlying substance of atoms were gradually revealing itself in a concrete and tangible form. Although scientists may regard the existence of atoms as their starting point, Cassirer, following Duhem and Ostwald, argues that a closer examination of the empirical grounds for their claims reveals precisely the opposite. One starts with a series of experimental laws (the law of the conservation of mass, definite, reciprocal and proportions) that reveal that, Unlike physical mixtures, chemical compounds have fixed compositions by mass. At this point, concept of \s{combining weight} naturally transforms into the concept of a discrete \s{atomic weigh}, as arising from whole-number combinations of discrete units: \qt{It becomes clear that the atom is never the given starting-point, but always only the goal of our scientific statements}{Das Atom ist, wie jetzt sogleich deutlich wird, niemals der gegebene Ausgangspunkt, sondern immer nur der Endpunkt unserer wissenschaftlichen Aussagen} \SF{276}{208}.


%\todo{Once the fixed atomic weight was established, the differing equivalents were explained by valence}
%\todo{the amount of an acid required to neutralize a specific base. If elements only ever formed a single compound with one another, these macroscopic equivalent weights and microscopic atomic weights would be effectively indistinguishable}
%\todo{Because atoms combine in simple whole-number ratios (1:1, 1:2, 2:3), the macroscopic weights measured in the laboratory are mathematically constrained to be simple multiples of the underlying microscopic weights.}
%\todo{The pure equivalent concept becomes ambiguous when elements form multiple compounds. If the equivalent is simply the combining weight, an element like carbon would have multiple, different equivalents depending on whether it was forming carbon monoxide or carbon dioxide. To maintain a consistent table of equivalents, Wollaston was forced to arbitrarily state that carbonic acid contained one equivalent of carbon and two equivalents of oxygen. In doing so, he silently introduced discrete multipliers, relying on the very concept of countable units that the equivalent was supposed to avoid.}
%\todo{Because he tried to identify the equivalent directly with the combining weight, a substance like carbon in carbonic acid forced him to assume two equivalents of oxygen for one equivalent of carbon. Consequently, the combining weight and the equivalent were no longer identical, introducing a significant confusion into the science that persisted for decades By assigning a definite, unalterable weight to the indivisible atom of every element, Dalton's theory provided the necessary, consistent structural explanation for these numerical regularities. It was Wollaston's subsequent attempt in 1814 to replace the Daltonian atom with the equivalent that actually retarded the general acceptance of Dalton's doctrine.}
%
%\todo{alency represents the discrete number of combinatorial capacities possessed by a single atom, which are pairwise saturated when molecules are assembled.}


%The wealth of content it acquires in the course of scientific research therefore never really concerns the atom itself, but refers instead to an empirical "subject" of a different kind. While we appear to investigate the atom itself in its manifold determinations and states, we are in fact establishing new relations among these various groups of states. We speak of the number of atoms contained in a given volume of a gaseous substance, but what we actually express is a relation that, according to Gay-Lussac's law, obtains between the numerical values of the densities of gases and those of their combining weights. We attribute the same heat capacity to the atoms of all elementary substances, thereby expressing the fact that, if we arrange the combining weights of the chemical elements in one series, a, a′, a″, …, aₙ, and the numerical values of their specific heats in another series, b, b′, b″, …, bₙ, there is a one-to-one correspondence between the two series, insofar as the products ab, a′b′, a″b″, and so forth, all have the same constant value.

% the mass of a gas per unit volume. Combining weight: the relative mass in which a chemical substance enters into combination with other substances.


%What appears to be the attribution of an identical heat capacity to the atoms of different elements thus expresses a quantitative relationship between two independently determined series. 

% on the other, the Dulong--Petit law establishes a quantitative relation between these same atomic weights and the specific heats of solid elements. 

%The same atomic weight can be established through chemical multple proportions, inferred from gas and vapor densities, or estimated from specific heats by means of the Dulong--Petit law. 

%Similarly, the chemical atom provides a conceptual representation of the simple whole-number ratios exhibited by the laws of chemical combination. Rather than deriving these regularities from the intrinsic properties of independently established microscopic substances, atomic theory expresses them through a unified system of quantitative relations.


%%**whole multiples** (*ganze Vielfache*) of a characteristic equivalence number, the concept of combining weight naturally transforms into the concept of a discrete **atomic weigh**

%Like Duhem or Ostawlad  Cassirer, seems ofacat the approacti Duhem's u the atom-conce is not a substac-concept\todo{improve}. 

Cassirer rejects Wurtz naive atomism about atom, however, he rejects Duhem's anti-atomism. For Cassirer, the alternative is misleading. The question lies not in the \qt{metaphysical claims about the existence of atoms}{metaphysischen Behauptungen über die Existenz der Atome} but in \qt{the logical function that the concept of the atom \textelp{} performs}{logische Funktion, die der Atombegriff \textelp{} besitzt} \SF{**}{**}. If one confines oneself strictly to observational data, macroscopic experience one might accept as a matter of fact that elements combine in several distinct ratios. Daltonian atomism, by contrast, forced chemistry to resolve the ambiguity by resorting to a combination of independent physical laws\footnote{Here Wallston's chemical equivalents and Dalton's atomic weights discused \citep[]{Ladenburg1902}}. While we appear to investigate the atom itself and its weight as a property, they instead used the sequence of atomic weights as a mediating series that makes it possible to coordinate otherwise unrelated domains of experience. The Avogadro--Cannizzaro method, for example, establishes a relation between the vapor densities of gaseous substances and their molecular weights, thereby providing an independent criterion for determining atomic weights:

\begin{equation*}
\rho_1, \rho_2, \ldots, \rho_n\,,
\end{equation*}
%
to their molecular weights, relative molecular weight of any gaseous substance is directly proportional to its gas density, from which the atomic weights of their constituent elements, 

\begin{equation*}
a_1, a_2, \ldots, a_n\,,
\end{equation*}
%
can be determined by taking chemical composition into account. On the other hand, the Dulong--Petit law establishes a correspondence between these same atomic weights and the specific heats of solid elements, 

\begin{equation*}
b_1, b_2, \ldots, b_n\,,
\end{equation*}
%
so that the product atomic weight and spefici weight is a fixed constant.
The same series of atomic weights $a_n$ thus connects two apparently unrelated empirical domains, the behavior of gases and the thermal properties of solids \SF{276--77}{208}. 

The question of the reaolty of atoms, this is here replaced by their succes of the atom-concept in organizing in establish a univocal corodatn beteen otherwise unrealtied series: \qt{The atom functions here as the conceived unitary center of a system of coördinates, in which we conceive all assertions concerning the various groups of chemical properties arranged}{Das Atom fungiert hier gleichsam als der gedachte einheitliche Mittelpunkt eines Koordinatensystems, dem wir alle Aussagen über die mannigfachen Gruppen chemischer Eigenschaften eingeordnet denken} \SF{277}{208}. The various and initially heterogeneous numerical sequences acquire a univocal coordination \origg{eindeutige Zuordnung} as we refer each of them back to a mediating sequence of atomic weights. The atom with a determinate atomic weight is not the underlying \s{bearer} of different properties that we discover; rather, the sequence of atomic weights is the \qt{common sequence of comparison}{gemeinsame Vergleichsreihe} mediating between different numerical sequences: \qt{In truth, we are concerned not so much with relating the diverse sequence \origins{Reihen} to the atom, as rather with \emph{relating them reciprocally} to each other through the mediation of the concept of the atom}{In Wahrheit handelt es sich darum, die verschiedenen Reihen nicht sowohl auf das Atom, als vielmehr, durch den Mittelbegriff des Atoms hindurch, wechselseitig auf einander zu beziehen} \SF{277}{208}. 

The empirical laws of multiple combination motivate the introduction of \emph{discrete} atomic units, while other laws make it possible to determine their relative weights and integrate them into a common system. The attempt to determine unambiguously \origg{eindeutig zu bestimmen} the atomic weight of each element requires drawing upon ever new empirical laws governing different domains of chemical and physical phenomena as criteria \SF{**}{**}. As this determination progresses, the range of empirical relations itself expands accordingly: \qt{If we imagine this progress brought to completion, then the totality of possible relations that the individual series can enter into with one another would be encompassed and expressed precisely in the \s{absolute} atomic weights}{Denken wir uns diesen Fortschritt vollendet, so wäre somit gerade in den »absoluten« Atomgewichten das Ganze der möglichen Verhältnisse, die die Sonderreihen unter sich eingehen können, befaßt und ausgedrückt}, or better the absolute value of the relative atomic weights \SF{**}{**}. On the other hand, by attributing observations concerning vapor density, heat capacity, isomorphism, and so forth to one and the same subject, we organize disparate known empirical laws into a system. On the other hand, in the attempt to determine the atomic weights, we discover new empirical laws that were previously unknown. Cassirer can thus conclude with a characteristic neo-Kantian move: the atom should not be conceived as the \s{substantial unity} from which scientific investigation begins, but rather as the \s{projected unity} toward which it strives: \qt{Thus, the chemical atom is an \s{idea} in the strict sense Kant gave to this term}{So ist das Atom der Chemie eine »Idee« in dem strengen Sinne, den Kant diesem Terminus gegeben hat} \SF{**}{**}.

\subsection{From Substitution to Valence}

%Atomic-weight determination establishes a characteristic value for each element, coordinating different empirical relations to determine its specific identity. Chemical types instead classify compounds already distinguished from one another, organizing them into families through substitution relations. The question shifts from “What determines this element?” to “How are these compounds systematically related?”

%for they describe not only singular molecular events or relations be-
%tween individual substances, but mostly generalized chemical relations between
%(open) substance classes

\todo{chemistry uses laws, but to determine the indiviau.}
%hus, the periodic table itself serves to defend Mendeleev’s stulate of the individuality of elements against those chemists or physicists who wanted to reduce the specificity of each element to a single “originary” proto-element.

%Wald, F.: 1899, ‘Was ist ein chemisches Individuum?’, Zeitschrift f¨ur physikalische Chemie
%28, 13–16

%ts general organizing principle allows each chemical element to occupy a determinate position within a systematic order. This position establishes relationships with other elements and constrains expectations about its properties.
%For example, Mendeleev could predict properties of undiscovered elements precisely because their positions within the system were already determined by the relationships among known elements.

%Particular bodies instantiate the laws.	Each element occupies a determinate position within the system.

%For Venel, chemists, unlike physicists, try to understand material substances in their complexity and their individuality
%The idea is not ... particualrlu

%Physical theory organizes laws. Chemical theory organizes substances.

From a logical point of view, one can say that the first stage in the development of chemistry was the ability to determine the relative atomic weight of elements considered individually. In Cassirer reconstruction, the second stage was to establish how they are systematically ordered in relation to one another, arranging them into classes of a distinctive character. Following his sources, Cassirer points out that \qt{[t]he empirical facts that first lead to such relative distinctions and groupings within the overall system are given in the relations of chemical \ls{substitution}}{Die empirischen Tatsachen, die zunächst zu derartigen relativen Unterscheidungen und Zusammenfassungen inner|halb des Gesamtsystems hinführen, sind in den Verhältnissen der chemischen Substitution gegeben} \SF{**}{**}. It was Jean-Baptiste Dumas who noted that for every volume of hydrogen removed from a compound in acetic acid by chlorine, an equal volume of chlorine is taken up producing trichloroacetic acid which reacted in almost identical ways to the original. This result challenged the Berzelius's dominant dualistic assumption that molecules must be split into positive and negative halves so that replacing a \s{positive} hydrogen atom with a \s{negative} chlorine atom was physically impossible \SF{**}{**}\footnote{On Berzelius see \citep[lect.~VI]{Ladenburg1902}}. \cop{Yet as the theory of substitution gains ground, the opposite view increasingly prevails: even entirely dissimilar substances can replace one another in certain compounds without altering the nature of the compound}. Examining how different elements substitute for one another reveals basic rules governing these substitutions: \qt{The form of substitution can be fixed and expressed once and for all by means of certain \emph{numerical values} that we assign to each element and that are now added to the number of its combining weight}{Die Form des Austausches läßt sich durch gewisse Zahlwerte, die wir jedem Element beilegen und die nunmehr zu der Zahl seines Verbindungsgewichts hinzutreten, ein für allemal festhalten und ausdrücken} \SF{**}{**}. 

%https://web.lemoyne.edu/~giunta/papers.html

To create a mathematical system, one needs a baseline of $1$. Just like hydrogen's mass is the baseline for equivalent weights\todo{check is baseline is mention}, hydrogen's substitution capacity is the absolute minimum combining capacity. If one wants to insert oxygen into a chemical framework, the cost of admission is the removal of exactly two hydrogen units. This establishes the relational ratios: $Cl \sim 1 H, O \sim 2 H, N \sim 3 H, C \sim 4 H$. This establishes a further basis for relating the individual elements to one another, a \qt{new characteristic constant}{neue charakteristische Konstante} for each elementary substance belongs to them independently of their chemical affinity with which substances does it tend to combine: \qt{The \s{valency} of the elements is the expression of a certain fundamental property in them, which belongs to them independently of their chemical affinity}{Die ,,Valenz" der Elemente ist der Ausdruck einer bestimmten Grund- beschaffenheit in ihnen, die ihnen unabhängig von ihrer chemischen Verwandtschaft zukommt} \SF{**}{**}. Without the concept of valence, one might merely record individual cases in which atoms of one element replace those of another in compounds. Valence brings them under a general rule, allowing one to determine which substitutions preserve a given pattern of bonding.

The \s{general rule} is that substitution must preserve the total valence: If hydrogen's base value is set to 1, oxygen's value is 2. Therefore, if you want to replace oxygen in a compound using hydrogen, you must supply twice the volume. Once again it was realized, that \cop{chemical substitution does not occur in arbitrary fractions; it operates \textit{per saltum}. One atom of Hydrogen is replaced by exactly $1$ atom of Chlorine, or $2$ atoms of Hydrogen are replaced by $1$ atom of Oxygen and so on}. Arranged according to this principle, chemical compounds fall into various  \s{principal types} \origg{Haupttypen} \SF{**}{**}. Compounds belonging to the same type can be derived from a common basic form through successive substitutions governed by the valences of the atoms or groups involved. Substitution thus becomes an instrument of classification. Methane gas provides a simple illustration. In this bulk substance, carbon exhibits a total combining capacity of four, which is entirely satisfied by a proportional mass of hydrogen. Because chlorine shares the exact same exchange rate (valence) as hydrogen, one can progressively react methane with chlorine to generate a series of new substances. Each successive substitution strictly preserves carbon's overall fourfold requirement: 

% it reveals a common structural pattern beneath differences in composition and physical properties. Methane gas provides a simple illustration. 
%Methane $\rightarrow$ Chloromethane $\rightarrow$ Dichloromethane $\rightarrow$ Chloroform $\rightarrow$ Carbon Tetrachloride.
%
\[
\mathrm{CH}_4
\;\longrightarrow\;
\mathrm{CH}_3\mathrm{Cl}
\;\longrightarrow\;
\mathrm{CH}_2\mathrm{Cl}_2
\;\longrightarrow\;
\mathrm{CHCl}_3
\;\longrightarrow\;
\mathrm{CCl}_4.
\]
%
Despite exhibiting drastically different physical properties in a flask, because each step involves trading exactly one equivalent of hydrogen for one equivalent of chlorine, these compounds are chemically understood as proportional modifications of the same fundamental \s{methane type}. The same classificatory principle is illustrated by the water type. In the water type, $\mathrm{H{-}O{-}H}$, divalent oxygen provides two bonding positions. Replacing one hydrogen atom with an appropriate monovalent organic group yields an alcohol, $\mathrm{R{-}O{-}H}$; replacing both yields an ether, $\mathrm{R{-}O{-}R'}$. Despite their different properties, water, alcohols, and ethers share the same basic pattern of two attachments to oxygen.

Following the logical rather than historical order, Cassirer could then treat the notion of valence as the fundament of the notion of \s{chemical type}. What interests Cassirer is the fact that \s{substitution} represents the exact moment chemistry abandoned traditional \s{classification} generic category-concepts \origg{{Gattungsbegriffe}} and mimicked mathematical \s{classification} based on function-concepts \origg{Funktionsbegriffe}:  

\qt{Here again, the concept of \s{type} comes into consideration not in its significance for the special problems of chemistry, but only as a paradigm for certain logical relations. Indeed, it exhibits with full clarity the characteristic fundamental features that the analysis of the exact scientific concept has already established in general. The chemical concept of type, too, is formed not on the model of generic concepts but on that of serial concepts. The various compounds subsumed under a principal type owe this grouping neither to an external similarity in their sensible properties nor to an immediate kinship in their chemical functions. They belong together insofar as they can be transformed into one another by virtue of the relation that holds between the valences of the individual atoms. The more remote members of the series, however, need not offer the nearer ones any \emph{analogy} other than the one established by this very law of derivation}{Der Begriff des \s{Typus} kommt hier wiederum nicht in seiner Bedeutung für die speziellen Aufgaben der Chemie, sondern nur als Paradigma für bestimmte logische Verhältnisse in Betracht. Er weist in der Tat in voller Deutlichkeit die charakteristischen Grundzüge auf, die die Analyse des exakten naturwissenschaftlichen Begriffs bereits im allgemeinen festgestellt hat. Auch der chemische Typenbegriff ist nicht nach dem Muster der Gattungsbegriffe, sondern nach dem der Reihenbegriffe gestaltet. Die verschiedenen Verbindungen, die unter einem Haupttypus befaßt werden, verdanken diese Zusammenstellung weder der äußeren Ähnlichkeit in ihren sinnlichen Eigenschaften, noch der unmittelbaren Verwandtschaft in ihren chemischen Funktionen. Sie gchören zusammen, sofern sie kraft der Bezichung, die zwischen der Valenz der einzelnen Atome besteht, incinander übergeführt werden können, wobei indes die entiernteren Glieder der Reihe zu den näheren keine andere Analogie mehr tarzubieten brauchen, als sie eben durch dieses Ableitungsgesetz selbst hergestellt wird.}[\SF{**}{**}]
%
Dumas still grouped compounds derived from one another by substitution that retained similar main chemical properties; Regnault extended the notion of substitution to encompass a whole sequence of reactions resulting in products that possessed completely different properties from the initial compound. The unity of the chemical type rests not on the similarity or analogy of the substances themselves, but on the rules governing their possible transformations through substitution governed by \s{valence}. Here the \emph{conceptual abstraction}, as conceived by traditional logic, establishes the unity of a class by identifying a characteristic common to its members, \emph{mathematical abstraction} establishes conceptual unity through their mutual substitutability \emph{salva veritate}. In geometry, two figures, however different they may appear at first sight, are regarded as equivalent if one can be obtained from the other through transformations belonging to the same group. Similarly, in chemistry, two compounds may be regarded as belonging to the same chemical type if one can be derived from the other, directly or through a succession of substitutions, even when their chemical functions and properties differ considerably.

%In Gauss's theory of congruences, for example, integers are said to be %\s{congruent modulo 5} when their difference is divisible by $5$. The integers $2,7,12,17,\ldots$ thus belong to the same equivalence class. Although these numbers are distinct, they can be substituted for one another in polynomial expressions without altering the result modulo $5$. Their unity is established not by abstracting common qualitative characteristics, but by a mathematical relation that determines their mutual substitutability. According to Cassirer, an analogous transformation occurred in chemistry. The unity of the chemical type rests not on the similarity or analogy of the substances themselves, but on the rules governing their possible transformations through substitution. Two compounds belong to the same chemical type if one can be derived from the other, immediately or mediately, through a series of substitutions, even when their chemical functions and properties differ entirely. 

Once again inverting the historical order, Cassirer concludes that the implications of this result become particularly evident when we turn to organic chemistry, as the theory of substitution acquired greater scope when chemists recognized that groups of atoms \origg{Atomgruppen}, as well as individual atoms, could function as substitutive units. The conversion of methanol into chloromethane, $$\ce{CH3OH + HCl -> CH3Cl + H2O}\,,$$ illustrates the persistence of a compound group through a chemical transformation. Here, \ce{CH3\bond{-}} functions as a compound radical \origg{zusammengesetztes Radikal} as Liebig defined it, while \ce{OH} is replaced by \ce{Cl}. A radical could replace one or more hydrogen atoms in a compound and was accordingly described as monovalent, divalent, and so on. For example, methyl (\ce{CH3\bond{-}}) has one point of attachment, whereas methylene (\ce{CH2\bond{-}}) has two, and so on.\footnote{Although these concepts were historically discovered in the reverse order, their logical dependency unfolds in precisely the sequence identified here, culminating in substitution.} Whereas Laurent considered radicals as \s{nuclei} possessing an independent existence in the compound, Gerhardt treated them as \s{residues} emphasized that their logical function of remain constant in a series of compounds \SF{**}{**}. In the expression \ce{CH3\bond{-}}, \ce{CH3} represents the radical, treated as constant across a series, while the empty place $\ce{\bond{-}}$ represents a replaceable atom or group. The empirical reality of \s{radicals} is not they correspond to \emph{actual} thing-like entities, but resolves into their logical function: they identify what remains invariant in a compound through possible substitutions, thereby allowing us to anticipate the possible behavior if substitution were performed suggests further possible members of the same series: \qt{The whole that we seek and toward which the concept is directed must not be conceived as an absolute being outside all possible experience; it is nothing other than the ordered totality of these possible experiences themselves}{Das Ganze, das wir suchen und auf welches der Begriff sich richtet, darf nicht im Sinne eines absoluten Seins außerhalb jeder möglichen Erfahrung gedacht werden; es ist nichts anderes als der geordnete Inbegriff dieser möglichen Erfahrungen selbst} \SF{**}{**}.  


%or Cassirer, the radical’s empirical significance lies in its role in organizing actual transformations and anticipating possible ones. It need not designate an independently existing little thing to perform that role

%In Cassirer’s reconstruction, chemistry was able to organize the manifold of substances, which was characterized by the assignment of a definite relative atomic weight to each element. 

%hey allow chemists to systematically classify substituti

\subsection{The Periodic System}

\todo{Duhem emphasises the important role played by the
prediction and actual synthesis of new substances in the development and consol-idation of the conception of structure based on substitution.}

Establishing a consistent scale of relative atomic weights was itself chemistry's substantial achievement. As Cassirer put it, each individual element thereby acquired, in Leibniz’s expression, a \s{characteristic number}, implicitly regarded as expressing its distinctive empirical character \SF{287}{**}. Whereas atomic weight distinguished the individual members, valence, an integer count of saturation units, provided a principle for classifying them according to their combining capacities, as distinguished from the affinity force of combination, measured by heat release. One might be tempted to explain these empirical facts by treating substances as made up of \cop{atoms as solid spheres of determinate weight endowed with physical \s{hooks}—if carbon had four hooks, they assumed it must always attach to four things}. However, a mechanical, erector-set model worked for early organic chemistry because carbon is famously consistent in how it bonds, but as inorganic chemistry advanced, it revealed that elements like nitrogen, phosphorus, and sulfur have variable valence. The real advance of notions like atomic weight, valence etc. is that the properties of as propertys of atoms as thinksl, into positons in the numerical series mediate the conversion of a qualitative concept into a position in a numerical series in which each element is defined by a pot \qt{This requirement now extends to all physical and chemical determinations recognized as dependent on definite numerical values \origins{Zahlwerten}{}}{o greift diese Forderung nunmehr auch auf alle die physikalischen und chemischen Bestimmungen über, die als abhängig von bestimmten Zahlwerten erkannt sind.} \SF{**}{**}. 





%The empirical sequence of qualitative elements is replaced by a sequence of possible elements conceived as members of a constructive sequence governed by a fixed rule.  To understand a phenomenon is to find something that plays the role of the general term of a sequence. From the substitution of an arbitrary natural number $n$ into this general member, one can calculate the individual member of the sequence that stands in the $n$th place. Once it is seen as the result of this calculation, the individual element appears as the expression of the entire sequence: 

\todo{arithmetrization}

%They, too, can no longer be conceived as a mere disorderly collection, but must be representable in their succession and gradual variation by an exact law.

\todo{improve}


%```
%Thus Locke, for example, develops his entire conception of both the tasks and the limits of natural-scientific inquiry by taking as his example the chemical knowledge of the elements and their properties. True knowledge, in his view, is attainable only where it is possible to gain universally valid insights into necessary connections. Only where all the properties of an object are perfectly intelligible and certain from its original nature, where it is thus possible to infer all its qualities directly from our acquaintance with the object and to determine them a priori, can one speak of knowledge proper, in the strict sense of the | word. |
%238	This requirement, however, which is met in all our "intuitive" judgments about mathematical relations, is one that our natural-scientific knowledge cannot satisfy. Here, where we are concerned merely with collecting and describing various facts of perception, it remains forever impossible to establish that dependence of the individual members on one another through which alone they could become a rationally connected and comprehended whole. However many properties of a substance we may discover through observation and experiment, the question of their inner connection is not thereby advanced a single step. No matter how many characteristics of gold we bring together, its ductility, its hardness, its resistance to fire, and so on, we cannot discover from them a single new determination; we can never understand the form of connection by virtue of which certain characteristics of one kind always correspond to certain others of another kind. An insight of this sort, which would first raise our knowledge of nature to genuine science, such as we find in mathematics, would be possible only if, instead of merely collecting observations about the empirical coexistence or the empirical incompatibility of characteristics, we could take up the problem "at the other end"; only if we could start from some determination of the essence of gold in order to derive deductively from it the totality of its secondary properties.[225] Modern science has partly fulfilled the ideal that Locke sketches here; but it first had to give this ideal itself a new meaning. It agrees with Locke that deriving the individual characteristics of a substance from its "substantial essence" exceeds the task of exact and empirical knowledge; but it does not therefore renounce every conceptual connection among the empirical data themselves. It gathers the multiplicity of the elements into a fundamental series whose members | follow one another according to a definite principle, and then determines the individual properties of bodies as functions of their position in this series. How the further characteristics follow from the assumed fundamental property, how a definite ductility and hardness, fusibility and volatility result from a definite atomic weight, admittedly remains unanswered here as well; nevertheless, the fact of this dependence itself is used in the attempt to calculate and predict certain particular | %239	qualities on the basis of certain general data. The functional connection established here therefore contains less, to be sure, than metaphysical insight into ultimate essential grounds; but at the same time it offers more than a merely empirical compilation of unconnected particulars. The order of the elements that now emerges offers at least an analogue of mathematics, and thus an analogue of exact and "intuitive" knowledge. We do not thereby penetrate any deeper into the absolute being of bodies; but we now grasp more sharply the rules of their systematic connectiont

%It gathers the multiplicity of the elements into a fundamental series whose members | follow one another according to a definite principle, and then determines the individual properties of bodies as functions of their position in this series.


%I kept the page markers (238, 239, |) and the note reference [225] where they appear in the source. For "Dehnbarkeit" I used "ductility," the term in English translations of Locke's gold example (Locke himself writes "malleableness"). "Feuerbeständigkeit" is rendered as "resistance to fire," Locke's "fixedness."


%Wurtz carefully separated **affinity** (chemical energy/force of combination, measured by heat release) from **atomicity** (saturation capacity or valence number)

%Wie es der eigentliche methodische Vorzug des Zahlgebiets ist, daß jedes Glied in ihm aus einer ursprünglichen Anfangssetzung nach bestimmten einheitlichen Regeln hergeleitet und konstruktiv entwickelt ist, so greift diese Forderung nunmehr auch auf alle die physikalischen und chemischen Bestimmungen über, die als abhängig von bestimmten Zahlwerten erkannt sind. Auch sie dürfen nicht länger als ein regelloses Beisammen gedacht werden, sondern | müssen in ihrer Abfolge und ihrer allmählichen Veränderung durch ein exaktes Gesetz darstellbar sein.




  %The \q{empirischen Realität} atomic weight and valence properties of substances in deterring non-arbitrary classification of substances, connecting atomic weight (the mass parameter) and valence (the substitution parameter) to predict chemical formulas, reaction pathways:  \q{Denn diese besagt nichts anderes als die durchgängige objektive Verknüpfung, in welcher ein individuelles »Ding« oder ein besonderes Ereignis mit dem Inbegriff der wirklichen und möglichen Erfahrungen überhaupt steh} \SF{**}{**}.

\todo{better adtion}

%\qt{This requirement now extends to all physical and chemical determinations recognized as dependent on definite numerical values \origins{Zahlwerten}{}}{o greift diese Forderung nunmehr auch auf alle die physikalischen und chemischen Bestimmungen über, die als abhängig von bestimmten Zahlwerten erkannt sind.} \SF{**}{**}.


Here the \qt{methodological advantage of the domain of numbers}{methodische Vorzug des Zahlgebiets} \origg{Zahlgebiet} becomes apparent \SF{**}{**}. Chemistry undergoes the same process of arithmetization that characterizes the history of physics, a process that culminates in the establishment of the \emph{periodic system}. Drawing on the more consistent atomic weights made available by Cannizzaro's 1858 reform, Cassirer attributed to Lothar Meyer the merit to have published his systematic table arranged the elements by increasing atomic weight and grouped their valence (combining capacity). In the 1884 edition of \textit{Die modernen Theorien der Chemie}, that Cassirer is referring to, \cop{Meyer arranged the elements in rows of increasing atomic weight from left to right, then aligned the rows so that elements with similar chemical properties fell into vertical columns, each representing a family of elements with the same \emph{highest possible} combining capacity or valence}. By arranging elements of equal valence in order of increasing atomic weight, Meyer revealed that the atomic weight differences ($\Delta A$) between analogous elements in a family follow step-like numerical progressions: approximately 16 for light elements, 46 for medium elements, and 88--92 for heavy elements. This table showed that chemical properties, specifically how elements bond with hydrogen or oxygen, repeat at regular \s{periodic} intervals.

Meyer plotted the physical atomic volume $V$ of the elements (calculated by dividing atomic weight by density in the solid state) against their atomic weight ($A$). Connecting these points produced a dramatic, wave-like periodic curve $V = f(A)$. Elements that occupy analogous positions on this curve, whether on its ascending or descending branch or at one of its maxima or minima, also prove to be analogous in their most important chemical and physical properties (valence, volatility, ductility, electrical and thermal conductivity, and so on). The high peaks of the curve were occupied by the light, low-density alkali metals (lithium, sodium, potassium, rubidium, cesium). The low troughs were occupied by dense, heavy transition metals (such as iron, cobalt, nickel, platinum, and osmium. Other physical properties—such as ductility, fusibility, volatility, thermal expansion, and electrical conductivity—increase and decrease in exact synchronization with this volume wave. As Cassirer emphasizes, \citet[176]{Meyer1884} himself observes that \s{matter} has now passed from the domain of the constants of natural science into the domain of the variables. Instead of regarding the characteristic properties of each element as independent \emph{constants}, chemists could now treat them as functions of a common \emph{variable}. This ordering relation is scientifically meaningful independently of whether atoms are understood as physically existing entities: \qt{The qualitative nature of the particular material is made mathematically conceivable by discovering a point of view from which the material can be arranged in sequence with a definite law of progression}{Das Qualitative der einzelnen Stoffe wird mathematisch faßbar und beherrschbar, indem ein Gesichtspunkt entdeckt wird, nach welchem es sich in Reihen von bestimmtem Fortschreitungsgesetz ordnet} \SF{289}{217}.  

%The \s{periodic system of elements'}had previously pointed, though with a somewhat vague gesture, toward a uniform \s{structural relations} of the various chemical atoms, a causal explation of such relatioos 

%functioned as **empirical constants** inferred from macroscopic mass ratios long before subatomic physics explained *why* matter is quantized


%\q{At the same time, however, this solution immediately leads to a new question. The problem arises of deriving the atomic weights, which were at first introduced only as discrete values, from one another by continuous variation, and of determining the law according to which the dependent properties would have to change under such a variation} \SF{**}{**}. %. Indeed, the organization of matter into numerical patterns, summed with the structural relations, into causal connections.



%The path of the development of chemistry is thus determined. The quantitative are replaced by a numerical series.  The numerical values obtained through observation are immediately arranged into series that proceed according to a rule, in which the subsequent members can therefore be determined from the preceding ones. 

Chemistry is now able to identify \cop{individual chemical elements as discrete arithmetical entities defined by their position in the periodic system}. At this stage, the relative atomic weight was assumed to be the ultimate, immutable, discrete constant \origg{starre Einzelwerte}. If atomic weights remain static, discrete constants, science is forced to accept them as unexplained empirical \s{givens}. Yet the constancy of relative atomic weights cannot be proved to be ultimate. Indeed, Cassirer could point to  the discovery of phenomena like radioactivity, which seems to suggest that the chemical atom might have its status as an ultimate, indivisible entity and \qt{dissolves into a system of electrons}{löst sich in ein System von Elektronen auf} \SF{**}{**}. All measurement methods led, with increasing accuracy, to the same values of the charge and mass of an electron. In this convergence, science might have found a new \qt{relative point of stability}{relativer Ruhepunkt} \SF{**}{**}. However, the new constant thus obtained might turn out to be a variable. \hide{Die neue Einheit, die auf diese Weise gewonnen wird, kann indessen selbst wieder nur als relative und somit im Prinzip veränderliche gedacht werden} \SF{**}{**}. Cassirer can therefore conclude that scientific inquiry never establish \s{atoms} as absolute elements of reality independent of any theory; the empirical reality of the \s{atom-concept} is the ideal limit of the series of changing theories about the atom: \qt{The convergence of the series \origg{Reihe} takes the place of the external criterion of reality}{Konvergenz der Reihe tritt an die Stelle des äußeren Maßstabes der Realität} \SF{**}{**}.  

\section{The Polemic with Rickert. The Case of Isomerism}
%
Just as physics replaces the classification of perceptible colors with a numerical ordering of monochromatic light according to wavelength, chemistry replaces the classification of substances by their qualitative properties with a numerical ordering of elements according to their relative atomic weights. However, Cassirer insists the chapter on chemistry was not simply as an appendix to previous sections on physics. Chemistry reveals an aspect of scientific concept formation that an exclusive focus on physics would have left unexplored. As Cassirer pointed more opening the next section, \qt{[t]he real \emph{methodological} interest of the construction of the concepts of chemistry lies in the fact that in it the \emph{relation of the universal to the particular} is set in a new light}{Das eigentlich methodische Interesse der chemischen Begriffsbildung liegt darin, daß in ihr das Verhältnis des Allgemeinen zum Besonderen in eine neue Beleuchtung gerückt wird} \SF{292}{221}. When we consider physics in isolation it seems that to generate the traditional antinomy: \s{knowledge} aims at the universal in the form of laws, while genuine \s{being} belongs to individual substances. Individual phenomena enter the picture chiefly as instances through which those laws are illustrated. From this perspective, one might reasonably conclude that increasing scientific generality entails a progressive neglect of individuality. According to Cassirer, in contrast to physics, chemistry seems from the beginning to try to grasp the qualitative features of individual things \origg{Einzelding} (such as gold, water, or specific organic compounds). 


Cassirer’s aim in his discussion of chemistry was to show that these two directions of concept formation in the natural sciences are mutually compatible and, indeed, inseparable. It is worth noting that the expression \s{concept-formation} \origg{Begriffbildung} does not play any role in the other works of Marburg school. For example, the term \german{Begriffsbildung} does not compare at all in Natorp \citeyear{Natorp1910a} book. On the contrary, prominence is given to this term by the so-called Heidelberg-Freiburg school. In particular Heinrich \citet{Rickert1902}\footnote{Windelband in Heidelberg and Rickert in Freiburg}  as famously to emphasis \s{limits of concept formation}. The more universal a scientific law or concept becomes, the further it moves from actual reality, creating an unbridgeable gulf between natural scientific concepts and actual empirical reality. To grasp reality in its uniqueness, human knowledge must adopt a different, complementary methodology—the individualizing method of the historical and cultural sciences. By framing his work in terms of the debate over scientific concept formation, Cassirer clearly sought to expose the inadequenci of Rickert's sharp distinction between the \emph{generality} of natural science and the \emph{individuality} of history. However, it is only after treating chemistry that Cassirer was in the position to launch a direct attack against Rickert in the concluding section \rom{9} of the first part of the book.

%By framing his work as contribution to the scientific concept formation, Cassirer clearly sought to enter in the debate opened by Rickert and expose the inadequacy of his position. This critique comes to a head in the concluding section \rom{9} of the book's first part. It is precisely chemistry that provides the basis for his challenge to Auch Rickert trennt durch einen scharfen Schnitt das Naturwissenschaftlich-Allgemeine vom Historisch-Individuellen ab. 

%The neo-Kantian split is his most telling example because the two schools share a starting point: both are Kantian and both reject psychologism

%Heinrich Rickert, Die Grenzen der naturwissenschaftlichen Begriffsbildung. Eine logische Einleitung in die historischen Wissenschaften, Tübingen/ Leipzig 1902, p. 255 [»[…] empirische Wirklichkeit […] wird Natur, wenn wir sie betrachten mit Rücksicht auf das Allgemeine, sie wird Geschichte, wenn wir sie betrachten mit Rücksicht auf das Besondere.«]


%This distinction became the basis of Rickert’s whole theory of historical knowledge. »Empirical reality becomes nature, if we consider it with regard to the universal; it becomes history, if we consider it with regard to the particular.«

%The pure »sciences of laws« are here contrasted with the »sciences of events,« the »nomothetic« procedure of natural science with the »ideographic« procedure of history. Rickert elaborates this distinction in his work,

%The ultimate goal toward which the physical sciences, like all other natural sciences, strive consists in removing empirical intuition from the content of their concepts. 

%Sigwart, Logik, Bd. I, S. 45 ff., 325

If Rickert’s reasoning is sound, Cassirer continues, scientific research has hitherto labored under a peculiar misconception about its own ultimate goal. Science seeks to know \s{reality}, yet the very instrument it employs, the natural-scientific concept, appears to frustrate this aim: \q{What was intended to sharpen our grasp of the particulars of empirical intuition must now be recognized as blunting it. To grasp reality conceptually is thus to destroy its most distinctive and fundamental content} \SF{**}{**}. According to Cassirer, however, this conclusion follows not from the nature of scientific concept-formation as such, but from the premises underlying Rickert's account of the latter. As one might expect, according to Cassirer, Rickert remains implicitly committed to the prevailing view of the universal concept is ultimately genus-concepts \emph{Gattungsbegriff} a representation of what is common and therefore as incapable of grasping the particular as particular. Following \citet{Sigwart1889} Rickert claims that concept-formation is nothing but an extension of word-formation:  \cop{Just as, for instance, no fully determinate intuitive content corresponds to the word \s{bird,} but rather only certain vague outlines of shape are retained in it together with the indistinct representation of the movement of flight, so that a child may also call a flying beetle or butterfly a bird |, so the same holds originally for all our general representations.}.  General scientific concepts also arise from indeterminate, faded representations that preserve only a few shared characteristics.Scientific concepts make classification more precise, but on this model they merely refine the same abstractive procedure, moving still further away from concrete intuition.

%Cassirer connects this position with Sigwart's psychological account of general representations: 


For Cassirer, however, exact scientific concept formation is not model on the refining the meanings of words, but it follows the logic of modern mathematics. The fundamental logical model of scientific thought is the sequence-concept \origg{Reihenbegriff}, whose prototype is the number sequence. Within the sequence of natural numbers, each number is a \emph{distinct} \s{individual}, uniquely determined by its position in the sequence. Yet all natural numbers are generated through the repeated application of one and the same rule, the successor function. A determinate number thus breaks with the traditional logical schema, which recognizes concepts only as general genus concepts subsuming a plurality of instances. \s{Two} and \s{four} do not exist as genera instantiated in all concrete pairs or groups of four objects of the same succession function. For Cassirer, numbers provide the paradigm of conceptual individuality. The transition from the individual to the general concept does not involve disregarding particular differences. Generality is thus achieved not through the progressive omission of differences, but by transformation of a constant into a variable. The individual thus appears not as something left behind by the general concept, but as a particular determination within it\footnote{Cassirer can draws on Rickert's student \citet{Lask1902} acknowledges that mathematical concepts possess a \s{concrete universality}: unlike ordinary general concepts, they can generate their individual instances through construction}. 
 
% Rather, it consists in treating specific determinations as variables that can assume different values within a systematically defined range which is generated according to a rule absriac tunviera is the \s{concrete universality}
 
The same logical principle underlies the formation of general physics concepts. Consider a particular quantity of gas. Initially, we might describe it in terms of its indefinitely many perceivable characteristics. Physics first replaces the real gas with an idealized model whose behavior is characterized by only a few variables: pressure $p$, volume $v$, and absolute temperature $T$. A possible state of the gas is represented by a triple $(p,v,T)$, whose components range over appropriate domains of possible numerical values ($\mid p>0, V>0, T>0\}$). The here mathematically possible does not yet mean physically admissible \SF{**}{**}. The choice of these state variables reflects the fact that their values are connected by a functional relation, expressed in the ideal gas law as  

\begin{equation*}
pv=RT\,,  
\end{equation*}
%
which restricts those candidate states. The ideal gas is thus an abstract representation of a real physical system: it describes how a gas \emph{would} behave under idealized conditions in which additional physical effects are neglected. In this model, all gases are treated as equivalent. Their substance-specific characteristics are disregarded, and their behavior is represented by the same functional relationship. At this point, Rickert's objection might still appear justified: \qt{scientific research cannot do without, alongside | the \s{given} elements of perception, purely ideal limiting concepts that cannot be exhibited in any direct experience; that it cannot do without the \s{unreal} alongside the \s{real.}}{daß die wissenschaftliche Forschung neben | 249	den »gegebenen« Elementen der Wahrnehmung die rein idealen und in keiner direkten Erfahrung aufzeigbaren Grenzbegriffe, daß sie neben dem »Wirklichen« das »Nichtwirkliche« nicht zu entbehren vermag}. \SF{**}{**}

Rickert might therefore argue that the formulation of universal laws requires the elimination of individual differences. According to Cassirer, however, scientific concept-formation proceeds in a very opposite direction, it does not remove individual differences but it transforms them into numerical variables. To account for the differences between \s{real} gases and the \s{ideal} gas, physics moves to a more comprehensive functional relation that encompasses the ideal gas law as a special case. For example, the van der Waals equation, 

\begin{equation*}
\left(p+\frac{a}{v^2}\right)(v-b)=RT\,,
\end{equation*}
%
introduces two additional empirically determined constants, $a$ and $b$, whose values differ from one gaseous substance to another\footnote{I borrow this example from a later writing \DI{**}{**}}. These constants represent, approximately, intermolecular attractions and the finite volume occupied by molecules, allowing the same mathematical framework to distinguish substances whose behavior the simpler model treated as equivalent. The ideal gas law is recovered by setting $a=b=0$. Thus, the introduction of additional parameters does not merely refine the description of a particular gas; it transforms the original law into a special case of a more comprehensive family of possible functional relations. The more general mathematical framework thereby permits a more differentiated empirical characterization of particular gaseous substances: \q{Auch von dieser Seite her zeigt es sich somit, daß dem wissenschaftlichen Begriff als solchem | die Feststellung des Einzelnen keineswegs versagt ist} \SF{**}{**}. The movement toward greater systematic generality is simultaneously a movement toward the more precise determination of the particular.


%More general, because the earlier law is recovered as a special case. More differentiated, because distinct parameter values describe different gases.

%Scientific knowledge advances toward the particular not by abandoning general laws, but by embedding them within more comprehensive functional relations that permit increasingly precise empirical determinations. What was originally a universal law becomes one particular case within a broader system of possibilities. 

%, represents the relationship between pressure $p$, molar volume $v$, and temperature $T$ through a general mathematical function. Each particular value of pressure, volume, or temperature is itself a uniquely determined numerical magnitude within the relevant numerical continuum. 



%Rather than classifying an already given collection of actual objects according to common characteristics, a functional concept establishes a law governing a systematically determined range of possible objects. Particular cases are obtained by specifying values of the variables entering into this law. As Weyl would later emphasize, functional concept formation consists not in disregarding individual differences but in making particular determinations variable within a constructively accessible range of possibilities. The general concept is thus not merely a class encompassing given individuals, but a principle governing their possible generation and differentiation.


%However, the it is clear his concern logical relation between the universal and the particular, but a similarity relation between the concrete individual thing given in sensation and the individual considered by physics\footnote{By Weyl: We cannot agree with Cassirer's remark that in this procedure the more general concept is the richer; for the properties of the individual ellipse depend, in addition to the general form of the equation, on the specific values of the coefficients}.

\hide{Denn die primäre Aufgabe der Begriffsbildung ist es nicht, wie die Logik unter dem Zwange einer jahrhundertealten Tradition zumeist angenommen hat, die Vorstellung zu immer größerer Allgemeinheit, sondern sie zu wachsender Bestimmtheit zu erheben. Sofern vom Begriff »Allgemeinheit« verlangt wird, so ist sie doch | 252	nicht Selbstzweck, sondern sie dient nur als Vehikel, um zum eigentlichen Ziel des Begriffs, zum Ziel der Bestimmtheit zu gelange.}

%The decisive point is that scientific progress does not necessarily introduce more possible objects. Rather, it introduces more variables or relations through which previously indistinguishable objects become distinguishable. hat has increased is the discriminating power of the conceptual representation.

%Atomic weights and stoichiometry: Determine the elemental composition of compounds. Valence: Establish constraints on the possible combinations of atoms. Structural formulas: Represent the different connectivity patterns compatible with a given composition and valence. Stereochemistry: Distinguish different spatial configurations of atoms having the same connectivity.

%Quantitative resolution: We measure an existing variable more precisely.
%For example, improving the measurement of an atomic weight from two to six decimal places.
%Conceptual resolution: We introduce an additional variable or a new kind of relation.

%The path of the development of chemistry is thus determined. The quantitative are replaced by a numerical series. The atomic interpretation held that relative atomic weights are whole-number multiples of hydrogen taken as a unit. The numerical values obtained through observation are immediately arranged into series that proceed according to a rule, in which the subsequent members can therefore be determined from the preceding ones. Yet the constancy of relative atomic weights cannot be proved to be ultimate. The relative atomic weight was assumed to be the ultimate, immutable, discrete constant. \origg{starre Einzelwerte} If atomic weights remain static, discrete constants, science is forced to accept them as unexplained empirical \s{givens}. The constant might become a variable again. Indeed, this process was beginning to take place in Cassirer’s time, when the rigid chemical atom lost its status as an ultimate, indivisible entity \q{löst sich in ein System von Elektronen auf} of fixed charge-to-mass ratio. In this new constant, science found a new \q{relativer Ruhepunkt}. However, there is not. Cassirer can therefore conclude that scientific inquiry does not establish \s{atoms} as absolute elements of reality independent of any theory; the reality of the atoms is the ideal limit of the manifold, changing theories about the atom: \qt{The convergence of the series takes the place of the external criterion of reality}{Konvergenz der Reihe tritt an die Stelle des äußeren Maßstabes der Realität}.

%An initial space of possibilities. A scientific theory provides a framework within which phenomena can be represented and distinguished.
%2. An enlargement of that space. New empirical discoveries reveal distinctions that the original framework cannot accommodate.
%3. The introduction of more general laws. The enlarged space is organized through new constraints that distinguish admissible possibilities and establish systematic relationships among them.

%For the broader historical development, therefore, chemical substance space might be the more appropriate expression.

%First, they constrain the space of possibilities. Empirical discoveries may force scientists to enlarge the space itself by introducing new dimensions of determination.

\todo{same for law}

%Insofar as universality is not an end in itself but serves merely as a means of attaining the concept's true goal, namely, determinacy. 

% What physics for  the dynamical laws \origg{kausalen Abhängigkeit} of phenonema, Chemistry. Here, however, from the domain \s{chemical} establishments, structural relations \origg{Strukturverhältnisse} of individual substances, that is formulas. Alson in this case, once a new individual difference emerges, chemistry moves to a more general concept to grasp it.  Here lies the particular \emph{methodological interest} of the latter and where Cassirer believes that his treatment of chemistry pays off. 

The example of the gas laws for the primary task of concept formation is not increasing indeterminacy but rather bringing it to increasing determinacy\todo{Refenrece}. Rickert might still object that even the \s{real gas} thus conceived is not a genuinely concrete \s{individual}. According to Cassirer, however, this objection loses much of its force once we move from physics \emph{laws} of physics to the \emph{formulas} of chemistry. To establish this point, Cassirer appeals to the phenomenon of \s{isomerism} \todo{Ostlwad Duen}. Before the recognition of isomerism, chemists generally assumed that when two distinct substances share the exact same empirical or \emph{compositional} formula

\begin{equation*}
\mathrm{C}_2 \mathrm{H}_6 \mathrm{O}\,,
\end{equation*}
%
they are \s{indistinguishable} at the level of elementary percentage. \cop{Yet cases were discovered in which allegedly compositionally identical substances, in analogous circumstances, give rise to different reactions, undergoing different substitutions}. If chemistry followed the concept-formation traditional logic, trying to grasp the individual difference between the two isomers would require adding more \emph{concrete} sensory descriptions (color, boiling point, odor). On the contrary, chemistry proceeds in the opposite direction, introduced more \emph{abstract}, \emph{constitutional}, or structural formulas, which takes into account \qt{how atoms are arranged alongside one another and mutually linked to form the unified structure of the molecule}{wie die Atome sich nebeneinanderlagern und sich wechselseitig zu einem einheitlichen Aufbau des Moleküls verknüp|fen} \SF{**}{**}:


%For example, ethanol and dimethyl ether have the same molecular formula, \(\mathrm{C_2H_6O}\), but different structural formulas:

\begin{equation*}
\chemfig{CH_3-O-CH_3} \quad \quad \chemfig{CH_3-CH_2-OH}
\end{equation*}
%
specifies how those atoms are bound together through valency units to form the molecular framework. \todo{stereochemistry} Thus, individual differences of isomers are grasped by introducing a moving beyond simple weight to include the spatial configuration of atoms:

\qt{Thus, for example, the chemical ``concept'' of a particular substance is given by its structural formula, in which it is conceived as a specific substance with its distinctive structure, while at the same time being assigned to the various chemical ``types'' and thereby placed in a definite relation to the totality of other substances. The ordinary chemical formula, which indicates only its general composition but not the manner in which the individual elements are arranged, is here enriched by a wealth of new relations. The general rule now at our disposal also allows us to trace how and according to what law the given substance is transformed into others: it encompasses not merely the form of its existence at a particular moment, but the totality of its possible spatial and temporal phases. The further chemical concept formation advances, the more sharply its capacity to distinguish particular substances becomes articulated. Substances that, from the standpoint of the less developed concept, had to be regarded as alike because they were ``isomeric'' become clearly differentiated from the standpoint of the more developed concept and are distinctly demarcated in their individual characteristics.}{So ist etwa der chemische »Begriff« eines bestimmten Körpers durch seine Konstitutionsformel gegeben, in der er als besonderer Stoff in seiner eigentümlichen Struktur aufgefaßt, zugleich aber den verschiedenen chemischen »Typen« eingeordnet und damit in ein bestimmtes Verhältnis zu der Gesamtheit der übrigen Körper gerückt wird. Die gewöhnliche chemische Formel, die nur im allgemeinen seine Zusammensetzung, nicht aber die Art des Aufbaus der | einzelnen Elemente wiedergibt, wird hier durch eine Fülle neuer Beziehungen bereichert. Die allgemeine Regel, die wir jetzt besitzen, gestattet zugleich zu verfolgen, in welcher Weise und nach welchem Gesetz der gegebene Stoff in andere übergeht: Sie schließt nicht nur die Form seines Daseins in einem bestimmten Moment, sondern die Allheit seiner möglichen räumlich-zeitlichen Phasen in sich. Je weiter die chemische Begriffsbildung fortgeht, um so schärfer prägt sich somit die Fähigkeit der Unterscheidung des Besonderen aus. Stoffe, die vom Standpunkt des unentwickelten Begriffs als gleichartig – weil als »isomer« – zu bezeichnen waren, treten vom Standpunkt des entwickelten Begriffs deutlich auseinander und grenzen sich bestimmt in ihrer Eigentümlichkeit ab.}[\SF{**}{**}]
%
But the processes were not concluded. Compounds were found to share the same structural formula while exhibiting opposite optical behavior\footnote{The so-called enantiomer; one enantiomer rotates plane-polarized light to the right, while the other rotates it by exactly the same angle to the left}. Optical isomerism thus revealed the limitations of two-dimensional structural formulas, just as constitutional isomerism had exposed the limitations of compositional formulas. Once again, the \emph{individual} differences between substances could be captured by introducing more \emph{general} forms of representation. Van 't Hoof replaced the two-dimensional structural formulas  with three-dimensional stereochemical formulas. Cassirer resisted the idea that these formulas represent the actual spatial arrangement of atoms, they encode the range of possible behaviors of compounds under different chemical conditions: \qt{the constitutional formula is, precisely thereby, also the genuine scientific expression of the empirical reality of the body: for this reality signifies nothing other than the thoroughgoing objective connection in which an individual \s{thing} or a particular event stands with the totality of actual and possible experiences in general}{Indem die Konstitutionsformel diese Verbindung herstellt, ist sie freilich eben dadurch zugleich der echte wissenschaftliche Ausdruck der empirischen Realität des Körpers: Denn diese besagt nichts anderes als die durchgängige objektive Verknüpfung, in welcher ein individuelles »Ding« oder ein besonderes Ereignis mit dem Inbegriff der wirklichen und möglichen Erfahrungen überhaupt steht} \SF{**}{**}

The case of chemistry shows that, \textit{pace} Rickert, conceptual progress in science does not consist in disregarding particular differences in favor of increasingly abstract generalities. On the contrary, the development of more comprehensive relational systems allows for the increasingly precise determination of particulars: \qt{The individual in its peculiarity is threatened only by the universality of the blurred generic image, while the universality of a definite law of relation confirms this peculiarity and makes it known on all sides}{Nur die Allgemeinheit des verschwommenen Gattungsbildes bedroht das Einzelne in seiner Eigenart, während die Allgemeinheit des bestimmten Relationsgesetzes diese Eigenart festigt und allseitig kenntlich macht} \SF{299}{225}. Rickert might object that an individual chemical compound is a repeatable chemical kind rather than an individual in the strict sense. However, this objection would amount to conceding that his real concern is not that scientific concepts fail to capture individuality \origg{Einzelheit}, but that they fail to capture the concrete \s{thinginess} \origg{Dinglichkeit} of the objects of common-sense experience \SF{**}{**}. These, however, are distinct problems. What counts as an \s{individual} is not given to theory as an external \s{thing}, but is determined within the theory itself. Different forms of inquiry---chemistry, biology, history\footnote{Cassirer draws on a recent essay on individual causality by \citet{Hessen1909} to argue that even historical causality contains an element of universality. To explain the outbreak of a war, for example, historians routinely invoke a heterogeneous set of psychological, political, economic, institutional, and geopolitical regularities. Yet whereas the natural scientist’s primary objective is to establish generally valid relations, the historian’s task is to determine which regularities are operative in a particular case and how they interact to produce a unique outcome. The universal and the individual are therefore not opposed even in historical inquiry. On the contrary, the more thoroughly the relevant regularities are understood, the more precisely the individual event can be explained}, and even everyday experience---employ different criteria for individuating their objects. In this sense, and only in this sense, the integration of these different forms of \s{individuality} can never be completed. As one might expect from Cassirer, it remains an infinite task: \qt{Here, individual reality indeed reveals its fundamental character of inexhaustibility. Yet it is precisely the distinctive virtue of genuine scientific relational concepts that they undertake this task despite its being, in principle, incapable of completion}{Hier bewährt in der Tat die »individuelle« Wirklichkeit den Grundcharakter der Unerschöpflichkeit. Aber es bildet zugleich den charakteristischen Vorzug der echten wissenschaftlichen Relationsbegriffe, daß sie diese Aufgabe trotz ihrer prinzipiellen Unabschließbarkeit in Angriff nehmen} \SF{**}{**}.




% So, in chemistry, each compound is defined by its structural formula; in biology, an individual species is identified through its position within a system of evolutionary relationships, determined by common ancestry and descent.


%\todo{Die »Strukturverhältnisse«, die neben den Gesetzen der kausalen}
 
 
\section{A  battle over the Substance-Concept. The Meyerson--Cassirer Debate}

In the summer of 1910, Cassirer asked his publisher to send Meyerson the galley proofs of \textit{Substanzbegriff und Funktionsbegriff}. The com proofs of Cassirer's monograph were delivered with considerable delay.  As Cassir As once Cassirer they could continue their \qt{battle over the substance-concept}{[ein] Kampf über den Substanzbegriff} \letterecndp{Cassirer}{Meyerson}{27}{8}{1910}[153], was  by the editorila and full their recepti pos.  When Meyerson received the complete proofs of \SFt, his review of the \textit{Erkenntnisproblem} was already in the course of publication in the \textit{Revue de métaphysique}. In a letter from October 1910, Meyerson wrote to Cassirer: \qt{I am now in possession of your \textit{Substanzbegriff} in its entirety, albeit without the cover page. This is, however, unessential, since I shall in any case have the book bound. The (very long) article on your \textit{Erkenntnisproblem} has for some time been lying corrected at the editorial office, but, as I was recently informed, it may not yet appear in the next issue, for typographical reasons!}{Ihren Substanzbegriff besitze ich jetzt vollständig, allerdings ohne Deckblatt, was aber unwesentlich ist, da ich das Buch jedenfalls einbinden lasse. Der (sehr lange) Artikel über Ihr Erkenntnisproblem liegt schon seit geraumer Zeit korrigiert in der Redaktion, wird aber, wie mir neulich mitgeteilt wurde, vielleicht noch nicht im nächsten Heft erscheinen, aus typographischen Gründen!} \letterecndp{Meyerson}{Cassirer}{21}{10}{1910}[157]. The publication of Meyerson's review took several more months, and the review did not appear until January 1911.

%Cassirer is making a particularly important claim about the role of numerical determination in scientific knowledge. Chemistry demonstrates that the natural sciences do not merely establish universal causal laws. They also investigate the specific structural relationships that distinguish particular objects. These relationships can be expressed numerically, and the resulting numbers can themselves be organized into systematic sequences.

%\citep{Meyerson1911}

%The aspect of Meyerson's review that, at first glance, is most striking is its length. 


\subsection{Meyerson's Review of the \EPt}

In nearly thirty pages, Meyerson surveys both the first and second volumes of the \textit{Erkenntnisproblem}, offering several comments and critical remarks on each section. The review begins with a generally positive assessment: \q{Cassirer's book---as he himself warns us---is only apparently a book of history. In its essence, it is a systematic book, a work dominated by a theory that the author seeks to support through the study of the philosophical and scientific evolution of the modern period. This prevalence of a unified point of view [\ldots] constitutes one of the book's most fascinating features} \citep[100]{Meyerson1911}. Cassirer's attempt to integrate the history of science with the philosophy of science was hardly new to French epistemologists \citep[101]{Meyerson1911}, but this approach had remained largely foreign to the German philosophical tradition, with the notable exceptions of Friedrich Albert \citet{Lange1873} and Kurd \citet{Lasswitz1890}. Only recently had this approach begun to gain renewed prominence in Germany, \q{under the sustained influence of Hermann Cohen's school}. In light of this achievement, the \textit{Erkenntnisproblem} stood as \q{the most important work produced to date by a disciple of the master of Marburg} \citep[101]{Meyerson1911}. Meyerson praises Cassirer's remarkable erudition, which extends beyond the history of philosophy to encompass a wide range of developments in the history of science. In Meyerson's view, Cassirer's book was both a significant innovation and a model for the historical study of scientific thought, an approach that was far less familiar to philosophers and historians of philosophy at the time than it is today\footnote{On this point, see also the lecture Meyerson delivered on January 25, 1911, at the École des Hautes Études Sociales  (Meyerson 2011, 195--197)}.

Nevertheless, Meyerson could not refrain from pointing out a glaring shortcoming in Cassirer's book: \q{The author has not, it seems to us, paid sufficient attention to the history of chemistry}. Cassirer's analysis could have drawn \q{great advantage} \citep[102]{Meyerson1911} from the contribution of chemists to the problem of knowledge. In this regard, Meyerson seems surprised that Cassirer did not take into account the four volumes of Hermann \citets{Kopp1869} history of chemistry which would have considerably enriched the historical background of the \textit{Erkenntnisproblem}.\footnote{Hermann Kopp had been Meyerson's professor in Heidelberg between 1878 and 1880. \citep[31]{Bensaude-Vincent2016a}} In Meyerson's view, Cassirer's omission of the history of chemistry appeared to be more than a mere oversight. Meyerson suspects that, in order to justify his systematic premises, Cassirer deliberately obscured the history of chemistry because it did not align with his functionalist account of scientific cognition. 

Meyerson concedes that Cassirer might object that chemists remained largely isolated from the broader intellectual movement until the second half of the eighteenth century. Boyle is the only exception, being both a great chemist and a renowned physicist, and Cassirer could reply that he had in fact devoted considerable attention to Boyle's work. However, in the second volume of the \textit{Erkenntnisproblem}, in fact, Cassirer discusses Boyle's contribution to epistemology, referring solely to his anti-Aristotelian work \textit{De ipsa natura}, ith having replaced the Renaissance organicist conception of nature redefining nature as a \q{system of rules} \citep[312--313]{Cassirer1907b}  However this interpretation, for Meyerson \citep[125]{Meyerson1911}, entails an unforgivable \q{mutilation} of Boyle's thought. Upon closer examination, Boyle was not simply the supporter of a lawful conception of nature; he was -\q{a strict atomist} \citep[122]{Meyerson1911} and a proponent of a \q{corpuscular} theory of matter (an expression coined by Boyle himself): \q{His treatises are filled with it, and a large part of his activity was devoted to demonstrating the proposition that all changes can be due only to the displacement of molecules or corpuscles}. Because Cassirer overlooks Boyle's corpuscularism, his interpretation appears to Meyerson \q{incomplete and therefore, in a certain sense, misleading} \citep[124]{Meyerson1911}. Boyle was not seeking mere abstract mathematical rules, but the true physical causes and ontological reality behind phenomena.

If Cassirer had considered Boyle from this point, he would recognize the functional \emph{legality} of nature, but also the centrality of the concept of \emph{substance}. It is true that Cassirer does not resort exclusively to laws, by means of that notorious \q{subterfuge} of Comte and Mach, of claiming that atomists behave in a purely instrumentalist way, as if atoms exist:  \q{Rather, he seems to arrive at a system related to the logical theories of the Marburg school, in which the particular form of physical hypotheses would be justified through the application of the concept of function} \citep[**]{Meyerson1911}. As further evidence of Cassirer's attempt to conceal the history of chemistry, Meyerson \citep[103]{Meyerson1911} notes that Lavoisier the father of quantitive revolution of chemistry is never mentioned in the \textit{Erkenntnisproblem}. This omission, like that of Boyle's corpuscularism, did not seem coincidental. In Meyerson's view, Lavoisier's contribution to chemistry stood in radical conflict with Cassirer's depiction of the history of modern science. Rather than dispensing with the concept of substance in favor of a mathematical understanding of chemical transformation, Lavoisier systematically relied on the notion of the \q{thing in itself} to challenge the theory of phlogiston \IR{202--207}{222--227}. Kant himself, Meyerson \citep[103]{Meyerson1911} insists, was strongly influenced by Lavoisier, but Cassirer failed to acknowledge it for the sake of his own epistemological position.

The contrast between Cassirer and Meyerson ultimately turns on their different assessments of the relationship between scientific knowledge and \q{naïve realism}. For Cassirer, the development of science involves a progressive departure from the assumptions of ordinary experience. However, in this way in a form of crypto-positivism, that \q{purely phenomenal reality} or a \q{world of relations without a support \origins{un monde de rapports sans support}}. But is entirely foreign to scientific practice. When science destroys an old concept of substance (such as phlogiston), it does not reduce reality to pure mathematical functions; rather, it immediately posits a new objective substance (such as oxygen or atoms) to explain the phenomena. In his view, Cassirer's determination to reconstruct the history of science in accordance with his philosophical commitments leads him to \qt{over-interpret}{trop interpréter} the historical evidence \citep[121]{Meyerson1911}. The consequences are both historiographical and epistemological. On the one hand, Cassirer's reconstruction marginalizes aspects of scientific practice that do not fit his systematic interpretation. On the other, it unwittingly come to share the \q{positivist} conception of science as a system of laws and functional relations, neglecting the persistent role of substantial entities in scientific explanation \IR{1--46}{17--62}. 

\subsection{The \emph{Erkenntnisproblem} Revisited}
%
One may reasonably assume that, had Meyerson been familiar with \textit{Substanzbegriff und Funktionsbegriff} when he reviewed the \textit{Erkenntnisproblem}, he would not have raised this objection in the same terms. As we have seen, Cassirer not only devoted an entire section to concept formation in chemistry but also accorded it a distinctive \q{methodological interest} which is precisely the opposite of what Meyerson see to suggest \SF{292}{221}. According to Cassirer, chemical concepts show that scientific concept formation need not erase the distinctive characteristics of particular substances \SF{299}{224}. On the contrary, as chemical concept formation advances, it enhances our ability to distinguish particulars. From Cassirer's point of view, Meyerson, like Rickerts, implicitly assumes that concepts are formed through generic abstraction, that is, by progressively discarding individuals' distinguishing features. The more universal a concept becomes, the further it moves from empirical reality, opening an unbridgeable gulf between the concepts of natural science and the reality they seek to describe. Whereas Rickert draws from this premise the methodological conclusions about the limit of scientific concept formation, Meyerson carries it to its ultimate metaphysical consequence; like Rickert an \s{irrational} residue, that is, a residue that is irrational not in an absolute sense, but only in relation to the intellect's \s{discursive} capacity\footnote{It seems plausible that Cassirer, after the meeting with Meyerson, decided to add a brief comment on Meyerson's philosophy to one of the final drafts of chapter IX. Meyerson did inded. lthough Meyerson is named only in a footnote, Cassirer is clearly referring to his recent revival of Eleatic metaphysics.}

%\footnote{ According to Meyerson, Cassirer writes, \q{[i]n order to understand reality by our mathematico-physical concepts, [\ldots] we must destroy it in its real nature, i.e., in its multiplicity and changeableness} \SF{430}{324}. Since this process would undermine the very possibility of knowledge, science arrives at a compromise: the manifold physical qualities of things must ultimately dissolve into a single undifferentiated matter, an \q{ether}, which is nothing but the hypostatization of \qt{empty space devoid of qualities}{leeren eigenschaftslosen Raumes} \SF{430}{324}. The points of geometrical space are indistinguishable in themselves; qualitative differences are reduced to different arrangements of identical constituents. Ultimately, the distinction between matter and space disappears, and matter itself is reduced to space. To be sure, the homogeneous space of geometry seems to preserve the minimal degree of diversity necessary for scientific thought to operate without collapsing into absolute nothingness. Yet this equilibrium is precarious. In the absence of any physical reference, there is no criterion for identifying a point in space, and reason undermines the very identity it seeks to establish: \q{The eternally homogeneous, motionless `sphere of Parmenides' constitutes the ultimate goal to which all natural science unconsciously approaches} \SF{430--431}{324}.}

%s we have seen, for Meyerson in order to \s{explain} reality we have eliminate differences, scientific knowledge of reality can be considered complete only when reality itself has been fully annihilated. Reality escapes this otherwise inevitable fate only because it resists as something \s{irrational} the efforts of rational explanation

For Cassirer, Meyerson's philosophical stance did not arise from his interest in the history of chemistry; rather, he could probably retorded the accusation that historiography of chemistry was shaped by his misleading conception of scientific concept formation. If Cassirer had neglected Boyle and Lavooiser, Meyerson was just as culpable to have downplayed the role of Richter and Dalton. Nevertheless, Cassirer took Meyerson's criticism seriously. There is evidence that Cassirer read Meyerson's review of the \textit{Erkenntnisproblem} shortly after its publication. In late February 1911, he sent Meyerson a brief letter of reply:
%
\qt{I have just read your review of the \textit{Erkenntnisproblem} in the library, and I would like to express to you my sincere thanks. Of all the reviews my book has received---most of which you are familiar with---yours is, in every respect, the most stimulating and helpful. For you entirely grasp the whole work from the fundamental systematic perspective in which it was written, and thus your comments consistently address those points that were, also for me, the most important and truly decisive ones. Of course, you already know that I cannot agree with you on everything, but your critique contains such a sharp and concise formulation of the contrast between us that it already contributes significantly to a mutual understanding}{Soeben habe ich in der Bibliothek Ihre Besprechung des `Erkenntnisproblems' gelesen und möchte Ihnen doch sogleich meinen herzlichen Dank sagen. Von allen Besprechungen, die mein Buch erhalten hat---Sie kennen sie ja zum grossen Teil---ist die Ihre mir in jeder Hinsicht die anregendste und förderndste. Denn Sie fassen das Ganze durchaus von der systematischen Grundabsicht aus, in der es geschrieben ist---und so treffen Ihre Bemerkungen stets diejenigen Punkte, die für mich selber die wichtigsten und eigentlich entscheidenden waren. Daß ich Ihnen nicht in allem zuzustimmen vermag, wissen Sie ja---aber Ihre Kritik enthält eine so scharfe und praegnante Formulierung des Gegensatzes zwischen uns beiden, daß schon damit für eine gegenseitige Verständigung viel gewonnen ist.}[\letterecndp{Cassirer}{Meyerson}{21}{2}{1911}[161]]
%
Unfortunately, Meyerson's review appeared too late to be considered for the revision of the first volume of the \textit{Erkenntnisproblem}, whose second edition, at the time Cassirer wrote this letter, had already been published. Nevertheless, Cassirer assured Meyerson that he would pay close attention to his comments in revising the second volume. Before proceeding with its publication, Cassirer also visited Meyerson in Paris at the end of March 1911, thereby taking the opportunity to discuss the more delicate issues of the \textit{Erkenntnisproblem} directly with him.\footnote{See \letterecnd{Cassirer}{Meyerson}{30}{3}{1911}[164]. Meyerson had also sought to arrange a meeting between Cassirer and Gaston Milhaud, who however was forced to withdraw at the last minute because of family circumstances \citep[644]{EMLF}}.

The second edition of the second volume of the \textit{Erkenntnisproblem} was published in the autumn of 1911. The extensive chapter devoted to the development \q{from Newton to Kant} includes an entirely new section on the \q{problem of matter} in modern chemistry \citep[2:428--441]{Cassirer1911a} addressing the shortcomings highlighted by Meyerson. A letter to Meyerson from October 1911 confirms that Cassirer introduced them precisely in response to his review: \qt{Today you are receiving the second volume of the \textit{Erkenntnisproblem} in its second edition. I would be very pleased if you could occasionally find the time to look through it and share with me your opinion. You will notice for yourself that I have made extensive use of your review, and that it has been extremely valuable to me: for example, I have now added a new section on the development of chemistry, and I have also taken into account some details to which you drew my attention.}{Sie erhalten heute den II. Band des Erkenntnisproblems in zweiter Auflage. Ich würde mich sehr freuen, wenn Sie gelegentlich Zeit fänden, ihn durchzusehen und mir über ihn Ihre Meinung zu sagen. Daß ich Ihre Recension vielfach benutzt habe und daß sie mir gute Dienste geleistet hat, werden Sie selbst bemerken: so ist z. B. jetzt ein Abschnitt über die Entwickelung der Chemie neu hinzugekommen und auch manche Einzelheit, auf die Sie hingewiesen haben, berücksichtigt worden.} \letterecndp{Cassirer}{Meyerson}{25}{10}{1911}[178].  

%In the \q{new section on the development of chemistry}, Cassirer attempts to reply to Meyerson's objections. To this end, he investigates both Boyle's corpuscularism and Lavoisier's chemical writings, and compares them with his own functionalist perspective. However, Cassirer also tried to integrate this material into his own \s{grand narrative}

\todo{grand narrative}

Cassirer opens the chapter by reiterating that, contrary to modern physics, chemistry dominated the \q{stage of the medieval view of nature}, in which it is always certain characteristics of bodies and groups of bodies that, in this conceptual language, are reinterpreted as thing-like entities. However, now the victory of the mechanical conception of nature gives the concept of \s{element} an entirely new meaning. Boyle's \q{corpuscular philosophy} recast the \q{occult qualities} in terms of \q{distinctly cognizable properties, such as magnitude, shape, and motion} \citep[2:430]{Cassirer1911a}. As a result of this shift, the internal constitution of matter was no longer conceived as an \q{unfathomable ground} \citep[2:426]{Cassirer1911a} inaccessible to human cognition, but rather as an object that could be measured and interpreted in mathematical terms. In this way, Cassirer while acknowledging Meyerson's objection, he could integrate it into his own conceptual framework. In replacing old substantial forms with the system of rules governing the motion of corpuscle, Boyle replace the the old metaphysical substances with the laws of motion:  \q{The \emph{material} concept of nature is thus subsumed \origg{aufgehoben} under the \emph{formal} one, nature as a thing under nature as the totality of rules}{Der materiale Begriff der Natur ist damit | in den for|malen, die Natur als Sache in die Natur als Inbegriff der Regeln aufgehoben} \citep[**]{Cassirer1911a}. As Meyerson, Cassirer emphasize that Boyle's approach played a little role in the early development of chemistry itself. John \citet{Freind1709}, for example, attempted to derive the whole of chemistry by the Newtonian method, seeking to introduce the concept of attraction into it and to make it fruitful for explanation \citep[**]{Cassirer1911a}. However, these attempts to establish a chemical dynamics remained isolated, as empirical knowledge was not yet sufficiently advanced. 



%In Cassirer's reading, Boyle redefined nature not as a separate material entity, but as a formal system of rules (Inbegriff der Regeln) governing atomic motions, shapes, and size

%Cassirer argues that Lavoisier completed this transition by converting the chemical concept of substance into a purely epistemological correlate of quantitative conservation1. By establishing exact numerical mass balances ($m_{\text{initial}} = m_{\text{final}}$) and empirical element definitions, Lavoisier abandoned metaphysical searches for qualitative "essences" (such as phlogiston) in favor of mathematical equations and functional dependencie

Eighteenth-century chemistry continued to be dominated by the theories of the \s{phlogistonists,} who retained a qualitative conception of natural processes. It is through Lavoisier's work that the \q{scientific concept of chemistry itself is transformed}. If  \SFt Lavoisier is mentioned only in passing, in the second edition of the \textit{Erkenntnisproblem}, however, he takes center stage.  The decisive innovation lay not in the discovery of new facts, since phlogiston theory had already recognized that combustion could involve an increase in weight, but in the new logical interpretation of those facts. Whereas the phlogistonists regarded quantitative measurements as superficial, Lavoisier made the quantitative relations between the initial and final states of chemical transformations a fundamental theoretical criterion for determining what counts as a \s{substance}. Once again, whereas Meyerson regarded Lavoisier as exemplifying the centrality of the substance concept, Cassirer could incorporate Lavoisier into his broader account of the transformation of that very concept through the conversion of chemical experience into exact quantitative determinations. Admittedly, Lavoisier achieved this result by adopting a strictly empiricist approach. Nevertheless, as Cassirer argues, \q{Lavoisier's achievement implicitly represents one of the most important steps toward the recognition of the logical principles of empirical knowledge itself: for here the concept of substance passes definitively from its ontological into its modern epistemological-critical meaning, being conceived as a pure correlative concept to the idea of quantitative conservation.}{Dennoch bedeutet Lavoisiers Leistung implizit einen der wichtigsten Schritte in der Richtung auf die Erkenntnis der logischen Grundsätze der Erfahrungserkenntnis selbst: Denn der Begriff der Substanz geht hier endgültig von seiner ontologischen in seine moderne erkenntniskritische Bedeutung über, indem er als reiner Korrelatbegriff zum Gedanken der quantitativen Erhaltung gefaßt wird} \citep[**]{Cassirer1911a}

%Boyle atoms Thus, Lavoisier introduced a purely \s{quantitative} distinction between \s{element} and \s{compound}: compounds can be split into products that each weigh less than the starting material, elements cannot.

The Meyerson–Cassirer \s{battle over the substance-concept} was fought with historical weapons, but the battlefield was ultimately a purely systematic one. Cassirer was ready to accommodate Meyerson’s objections on historical grounds, though he entirely rejected their epistemological implications. Meyerson contends that scientists (such as Boyle with his rigid corpuscular atomism and Lavoisier with his conservation of matter) continuously appeal to something that remains constant beyond the apparent variety of the phenomena. By contrast, Cassirer could easily reinterpret the same episodes as confirming his narrative of the progressive reduction of the substance-concept to the function-concept (Boyle's atom moves according to mathematical laws, and Lavoisier )  Regardless of this divergence, in November 1911, Meyerson informed Cassirer that he found the second edition of the \textit{Erkenntnisproblem} \qt{considerably improved}{ganz wesentlich verbessert} \letterecndp{Meyerson}{Cassirer}{15}{11}{1911}[180]. Meyerson most likely appreciated Cassirer’s effort to engage with sources he had not considered in the first edition, even though Cassirer interpreted them within a narrative diametrically opposed to his own. If the Meyerson–Cassirer debate teaches us anything, it is perhaps the limitations of such \french{grand récit} that are either too narrow to accommodate the complexities of historical detail or, as in this case, so broad that virtually any historical evidence can be made to fit them. 

%Although both Meyerson and Cassirer claimed that their respective narratives emerged from the history of science, both ultimately too often yielded to the temptation to bend the history of science to fit their own narratives.

%The debate boroatb  that limits of the big Meyerson and Cassirer thus approached chemistry, but at the same time, the limit of integrating this material into a grand narrative became apparent. Although it imposed some constraints, the substance concept functioned as the removal of matter into a universal sphere of parameters. Both narratives are somewhat tentative, but also easily manipulated by choice and their interpretation. 



%James A. Secord, “Introduction” to the special issue “The Big Picture” (1993), British Journal for the History of Science 26(4): 387–389. Soraya de Chadarevian, “Microstudies versus Big Picture Accounts?” (2009), Studies in History and Philosophy of Biological and Biomedical Sciences 40(1): 13–19. Moti Mizrahi, “The Case Study Method in Philosophy of Science: An Empirical Study” (2020), Perspectives on Science 28(1): 63–88. Andrew Barry et al., “Conversation: The History of Science and the ‘Big Picture’” (2024), BJHS Themes 9: 225–235.


\section{Conclusion}
%
After 1911, the correspondence between Cassirer and Meyerson became more sporadic. However, in December 1912, just after the publication of the second edition of \IRt \citep{Meyerson1912}, in which he introduced multiple references to Cassirer’s works, Meyerson sent him two issues of the Bulletin de la \textit{Société française de philosophie} devoted to the discussion of Perrin’s study on Brownian motion. The brief letter accompanying the issues reads as follows: 

\q{At the same time, allow me to send you two issues of the Bulletin de la \textit{Société française de philosophie}, which you probably do not see regularly. They contain Perrin’s lecture on the new measurements of atoms, as well as the ensuing discussion. I also took part in the discussion and, of course, I attempted to turn the situation to my own advantage. Although you will probably disagree with almost everything, you might perhaps be interested in certain aspects of the whole subject. Should you write anything about it (for which I would be very grateful), I kindly ask you to send me a copy}[\letterecndp{Meyerson}{Cassirer}{1}{12}{1912}[195]]
%
Up to \citets{Perrin1910} work, the amounts of that element entering different molecules occur in integer multiples of its relative atomic weight. These relations leave the scale individual atomic masses undetermined \posscite[268]{Perrin1910}, since they do not establish Avogadro constant \(N_A\)---ratio between an amount of substance and the number of particles that it contains. Perrin showed that Brownian motion, electrical measurements, radioactivity, and other phenomena yield mutually compatible determinations of \(N_A\). Perrin takes the fact that this measurements \qt{converge precisely on the same number}{convergent ainsi précisément vers le même nombre} \citep[281]{Perrin1910} this as evidence for the actual existence of the atom each endowed with absolute weight.  \cop{Thusc the Brownian motion of small suspended particles demonstrates directly to our senses the presence of a molecular thermic motion}.

\todo{phenomeiast was antisitpisa, not wy tit was antirealistl that atom concet is  the exitence of atim, but tneht ature of the concept of atim }


For Meyerson, Perrin’s result showed that, despite the prestige of positivist ideas, science could not avoid reintroducing atoms as microscopic \s{things,} beyond the functional relations among measurable quantities. Cassirer’s response to Meyerson's letter reveals a cordial yet detached tone: \q{Thank you very much for sending me the report on the discussion at the \textit{Société française de philosophie}; I read it with great interest.} (**) However, he did not add further comments. Indeed, after the publication of \SFt, Cassirer moved away from \s{philosophy of science} until he wrote his relativity book a decade later \citep{Cassirer1921}. Here, Cassirer, only add sporadic to the role of atoms in chemistry, However, his stance seems to have been remained unchanged. The importance of chemistry in the history of the concept that constancy of atoms as thing, becomes the constancy in the constancy of a numerical value: \qt{fundamentally, such a realization, such a strictly quantitative meaning of the concept of the atom, was only reached in the beginnings of modern chemistry in the law of multiple proportion}{Im Grunde wird eine solche Erfüllung, wird ein streng exakter quantitativer Sinn des Atombegriffs erst mit den Anfängen der modernen Chemie im Gesetz der multiplen Proportionen erreicht.} \ER{**}{**}. As we have seen, according to Cassirer, with the further help of Avogadro’s Law, Dulong and Petit’s Law, and Mitscherlich’s Law, chemistry could remove the ambiguity in Dalton's law by progressively and univocally determining the \emph{relative} weights of atoms with respect to other atoms. Perrin's result could be seen as part of this process of further disambiguation. By fixing the exact scaling factor, Avogadro's constant, it becomes possible to univocally determine the \emph{absolute} weight of an individual atom in grams. In Cassirer's view the \s{reality} of the chemical atom is nothing but this convergence of progressively more stringent numerical determinations:  \qt{the atom shows itself thereby to be, not an absolute minimum of being, but a relative minimum of measure}{Das Atom erweist sich eben darin nicht als ein absolutes Minimum des Seins, sondern als ein relatives Minimum des Maßes} \ER{18}{360}. 

%Contray to Meyerson's was to denu, buyt to question of the atoms' existence, but with the logical function of the atom-concept \origg{Atombegriff} within chemical theory. In the historical development of scientific knowledge, the thing-concept \origf{Dingbegriff} of the early atomists is not a purely measure-concept \origg{Maßbegriff}. 

%The discoery of isotopes are the exact inverse of this indistinguishability, they have different atomic weights but the exact same chemical properties (e.g., Carbon-12 and Carbon-14). Because isobars are entirely different elements, they have different numbers of protons, even if their total mass (protons + neutrons) adds up to the same number.


% \todo{The atomic mass of an element is the weighted average mass of the atoms in a naturally occurring sample of the element.}. Once again, the atomic weight constant turned out to be a variable as a function of a constant of higher order, the atomic number 

%adding dimensions to the conceptual representation

% (van den Broek 1913)

Altug, share some anti-realist Cassirer do deny that atoms exitst, bit that atoms exit, but nature  atomc-oncet Cassirer points out that the absoluteness of this \s{minimum} always provisional to a certain stage of the development of science and cannot exist. Durch die Erkenntnis der radioaktiven Verschiebungs|sätze discovery that \s{isotopes} of the same element have different masses. The atomic weith gives average mass per atom when \(M\) is the molar mass of a mixture of differnt isotes \todo{Kramer}. Their Individual identity is not univocal determ atomic mass, Atomic weight could no longer be retained as the true ordering principle of the natural system, it proved \q{unexplained function of the true “atomic number.”} \citep{Cassirer1936a}. This true atomic number based on Bohr's atomic model, in a neutral atom, the number of electrons equals the number of protons \PSF{516}{442--43}. could the show that. Each element is again a position in a sequence, but now according to a deeper sequence principle \emph{Reihenprinzip}: \qt{This principle itself is never simply \s{given} \origins{gegeben} but represents a perpetual problem \origins{aufgegeben}, and one of the essential achievements of all natural science is to be sought in a steady advance toward a more complete solution of this problem}{Dieses Prinzip selber ist uns niemals schlechthin »gegeben«, wohl aber ist es beständig »aufgegeben« – und in der immer vollkommeneren Lösung dieser Aufgabe besteht eine der wesentlichen Leistungen aller NaturTheorie.} \PSF{518}{444}.\todo{Bohr atom}

%Das Bohrsche Atommodell erlaubte erst, im Bereich der Atomphysik gewissermaßen festen Fuß zu fassen. Es war der »Archimedische Punkt«, dessen sie notwendig bedurfte. Eine außerordent|liche Konzentration des bisherigen Wissens war er|reicht; 
%162	die Ableitung der Balmerschen Gleichung, der Rydbergschen Konstante, des Ritzschen Kombinationsprinzips war gegeben. Der Aufbau des periodischen Systems der Elemente gelang jetzt in der einfachsten und übersichtlichsten Weise, indem die Elektronenzahl des Atoms gleich der Atomnummer, d. h. gleich der Stellenzahl des Elements im natürlichen System, gesetzt wurde. Und dennoch blieb in dieser ersten Fassung der Bohrschen Lehre noch ein Bestandteil zurück, der theoretisch nicht völlig geklärt war und der insbesondere zu erkenntniskritischen Zweifeln Anlaß geben mußte.

%Das Atomgewicht konnte nicht länger als das wahre ordnende Prinzip im natürlichen System festgehalten werden; es erwies sich, daß es nur eine verwickelte und noch ungeklärte Funktion der wahren »Ordnungszahl« darstellte. Als diese wahre Ordnungszahl wurde die Anzahl der positiven Ladungseinheiten des Kerns oder die Elektronenzahl des Atoms gefunden

%They differ in their number of neutrons.

%substances with different nuclear masses (atomic
%weights) and different radioactive properties (if there
%are any), but with the same nuclear charge, the same
%number of electrons and, consequently, occupying the
%same place in the periodic system.

%(Atombau und Spektrallinien, S. 167).


%Das Atomgewicht konnte nicht länger als das wahre ordnende Prinzip im natürlichen System festgehalten werden; es erwies sich, daß es nur eine verwickelte und noch ungeklärte Funktion der wahren »Ordnungszahl« darstellte.


%Nicht nur ließen sich jetzt fast alle physikalischen und chemischen Eigenschaften eines gewissen Elements durch seine Stellung im System bestimmen – es zeigte sich auch, daß es auf Grund dieses Systems möglich war, die Existenz bestimmter noch nicht beobachteter Grundstoffe vorauszusehen und diese Voraussage durch die Erfahrung zu bestätigen

%Das Atomgewicht konnte nicht länger als das wahre ordnende Prinzip im natürlichen System festgehalten werden; es erwies sich, daß es nur eine verwickelte und noch ungeklärte Funktion der wahren »Ordnungszahl« darstellte. Als diese wahre Ordnungszahl wurde die Anzahl der positiven Ladungseinheiten des Kerns oder die Elektronenzahl des Atoms gefunden.[32] Damit war in der Tat alles erreicht, was innerhalb einer empirischen Theorie erreichbar ist und was eine solche sich als Ziel ihrer Begriffsbildung stecken kann



% of the element, which determines its position in the periodic table and was now identified with its nuclear charge, \todo{ All atoms of an element have the same number of protons, and every element has a different number of protons in its atoms} \PSF{516}{43}.


%Kern umgebenden Elektronen an (van den Brock 1913). Als die gesuchte physikalische Bedeutung der »Ordnungszahl« ergibt sich somit die Annahme, die die Grundlage der Bohrschen Atomtheorie bildet, daß die positive elektrische Ladung des Atomkerns von Element zu Element je um eine Einheit zunimmt. Der Atombau wird auf elektrischem Wege


%Aston found that also the masses of the atomic nuclei are, at least to a considerable approximation, integral multiples of the mass of the proton. If this did not become apparent in the atomic weights obtained by chemistry, the reason is to be seen in the fact that different atomic structures may belong to the same order number ('isotopes'), and that what chemists used to consider as a pure element frequently turns out to be a mixture of isotopes of different atomic weights; for these cannot be segregated by ordinary chemical means.

%The composition of the material world out of one or a few units, exist-ing in a huge number of completely alike specimens, must surely be looked upon as one of the most fundamental features in the nature of the uni- verse; and one that is most profoundly in need of interpretation.


%The various chemical elements differ from one another by the charges of their nuclei, which are integral multiples ne of e, where -e denotes the charge of the electron. If the atom is in an electrically neutral state ('non-ionized'), the factor n, called the order of the atom, coincides with the number of electrons revolving around the nucleus. There is no gap in the sequence of the orders of the various elements: n = 1, hydrogen; n = 2, helium; n = 3, lithium; . . . . Elements may change into one


%he old dream of the unity of all matter had certainly come a great step nearer to fulfillment by the discovery of the electron. But the positively charged nucleus of the atom about which the negatively charged electrons revolve like the planets around the sun still seemed to be a particle with a constitution of its own for each individual chemical element. It is a priori clear that beside the negative electron at least one positive brick is needed for the construction of all atoms. The execution of this idea of building up matter from two ultimate elementary units, the electron and the proton, presupposes of course that they are able to enter into a variety of combines that are held together by strong forces and react outwardly like solid atomic spheres. Hence sheer substance without force would never do.

%\todo{Onserivn taht diepen of maec, phenoe, and realit, mcirscorstrucural redifnto }. 

%As Meyrson \citet{Meyerson1931,Meyerson1933} despite their lip service to phenomenalism, could not avoid being \s{realist} about atoms, protons, and electrons \etc. 

%Auch die außerordentliche Fruchtbarkeit der Bohrschen Atomtheorie erschien vor allem darin gegründet, daß hier eine außerordentliche Fülle empirischer Einzeltatsachen gleichsam in einen Blick vereint und durch ein einfaches Modell von höchster anschaulicher Klarheit dargestellt war. Nichtsdestoweniger läßt sich gerade an dem Unterschied des Bohrschen Atommodells von der Vorstellungsweise, die in der »klassischen« Form der Atomistik herrschend war, der Wandel in der allgemeinen methodischen Grundansicht deutlich machen. 

%Tdie außerordentliche Fruchtbarkeit der Bohrschen Atomtheorie erschien vor allem darin gegründet, daß hier eine außerordentliche Fülle empirischer Einzeltatsachen gleichsam in einen Blick vereint 

%Die Antike und die Entstehung der exakten Wissenschaft[1] (1932) 


Bohr’s atomic model made it possible, for the first time, to gain a firm foothold, so to speak, in atomic physics,  derivations of Balmer’s equation, the Rydberg constant, Ritz’s combination principle, as well the periodic system. Here any doupt on the exitecen of atoms had t removed. However, for Cassirer, the distinction between \s{phenomenalism} and \s{realism} was once again the wrong fault line: \qt{We are concerned not so much with the existence of things as with the objective validity of relations; and all our knowledge of atoms can be led back to, and depends on this validity}{Es handelt sich nicht sowohl um die Existenz von Dingen als vielmehr um die objektive Gültigkeit von Relationen – und all unser Wissen von den Atomen läßt sich immer auf diese Gültigkeit zurückführen und geht in ihr auf} \DI{178}{143}. As one migth expect from a Kantian, the empirical reality of atomic \textit{entities} resolved completely objective validity of the atomic \textit{theory}. However, for Cassirer this \s{objective validity} is never achieved as ultimate, but it is always relative to a certain stage of the history of science. The succession of physical theories attribute to the atom a progressively more determinate form: \qt{But this \s{concrescence} \origins{Konkreszenz} of determining factors stands in an inverse relation to what we commonly designate as visualizable concretion \origins{Konkretion}}{Aber dieses Zusammenwachsen, diese »Konkreszenz« der Bestimmungsstücke steht zu dem, was wir unter anschaulicher Konkretion zu verstehen pflegen, im umgekehrten Verhältnis}  \DI{178}{143}. However, despite the loss of intuitive, pictorial models, physical \s{objectivity} is thereby increased rather than diminished, for \qt{true physical objectivity must never be sought in the direction of \s{naive} things, but rather in the direction of the critical concept of law \origins{Gesetzesbegriff}}{ie wahre physikalische »Objektivität » aber ist dadurch nicht vermindert, sondern vermehrt; denn sie kann von Anfang an nicht in der Richtung auf den 'naiven' Dingbegriff, sondern sie muss in der Richtung auf den kritischen Gesetzesbegriff hin gesucht werden.} \DI{178}{144\tm}.

If Cassirer borrw some anti-atomistic arugem, his concern is not the atocm realim atoms, but the fact that atom-concet is not a sustacen-concet. The entire history of the atom-concept represent for Cassirer a confirmation of this view. Here one observes indeed \qt{\emph{consolidation} of the atom-concept}{K o n s o l i d i er u n g des Atombegriffs}, but not an \qt{\emph{hypostatization}, its immediate reification}{H y p o s t a s i e r u n g, seiner unmittelbaren Verdinglichung ist} \DI{171}{144\tm}. In the sections dedicated to the history of the chemical atom, Cassirer had now clearly integrated some of Meyerson's suggestions, pointing out the role of Boyle’s corpuscularism is part of the hitory of quantitative chemistry fro the outset. However, that the Dalton's atomism chemical atom aquate aemprical realti of multy emrical alws, conserfavtion of matter to multimep pproptions: 

\qt{[T]he atom is here defined strictly substantially; but all statements concerning its nature \origg{Was} ultimately merge into statements concerning quantity-definite numerical relations \origg{Wieviel}. All the basic concepts and theories of chemistry, the concepts of valence and types, the concept of the radical as well as the theory of compound radicals all point in this direction. The conceptual language of atomism in conjunction with the general symbolic language of mathematics now permits that simple and magnificent development in the language of chemistry that is evident in the progress from simple to structural formulas. Stereochemistry, as established by van't Hoff, permitted the elucidation, through a complex construction of formulas, of an entirely new group of facts, the phenomena of stereoisomerism; it threw new light on the distribution of atoms in space. Thus the visual picture of the atom is modified and refined, but all these changes serve only the one purpose of circumscribing and representing an increasingly large domain of functional relations}{Wie bei Huyghens, so wird hier das Atom streng substantiell gefaßt; aber alle Aussagen über sein »Was« gehen zuletzt in Aussagen über ein »Wieviel«, in Aussagen über bestimmte numerische Verhältnisse über. Alle Grundbegriffe und alle Theorien der wissenschaftlichen Chemie – der Valenzbegriff und der Typenbegriff wie der Begriff des Radikals und die Theorie der zusammengesetzten Radikale – weisen in diese Richtung. Die Begriffssprache der Atomistik erlaubt jetzt, in Gemeinschaft mit der allgemeinen Zeichensprache der Mathematik, jenen einfachen und großartigen Aufbau der chemischen Formelsprache, wie er sich im Fortgang von den einfachen Formeln zu den Strukturformeln vollzieht. Die Stereochemie, wie sie von van’t Hoff begründet wird, erlaubt auf Grund einer komplexeren Gestalt der Formeln eine neue Gruppe von Tatsachen, insbesondere die Erscheinungen der Isomerie, zum Ausdruck zu bringen; sie gibt einen neuen Aufschluß über die Lagerung der Atome im Raume. So wandelt und verfeinert sich das anschauliche Bild des Atoms; aber alle diese Wandlungen dienen nur dem einen Zweck, einen immer größeren Kreis funktioneller Beziehungen zu umspannen und darzustellen}[\DI{183--84}{147}]
%
%The atom is ultimtly and hitorical construct,  %The mere fact that the century has seen a succession of different theories is not, by itself, grounds for pessimism or optimism. What would properly raise our suspicions is the opposite: a vigorous program of investigation into nature in which later researchers find no occasion to correct their predecessors. Our optimism or pessimism should rely on our examination of the details of the changes in electron theories. If these theories were to form a sequence of disconnected portraits, each merely answering to the transient expedients of the moment and each eradicating the content and successes of the earlier theories, then we could be excused for suspecting that we have just replaced one error with another as we pass from one theory to the next. But we do not have such a sequence. We have good reason to see our sequence of theories as correcting errors of former members while preserving their successes and providing richer and improved representations. We have shown that we can discern a
%
The discovery of the electron turned the chemical atom of Dalton into a composite object (Thomson's \s{plum pudding} model) and the discovery of the nucleus fixed into a spatially structured object (Rutherford's nuclear atom and Bohr's planetary model). From this point of view, as Cassirer put it, the history of the atom-concept appears a \q{kaleidoscopic} variety, a strange \q{phantasmagoria} of  disconnected portraits. However, the history has good reason to see our sequence as correcting errors of former members while preserving their successes \qt{in a single direction}{eine einheitliche \ls{Richtung}} \DI{187}{150}. It is this accumulation of historically stable properties that we call the \s{reality} of the atom: \qt{Thus we need not find an answer to the question of what the atom actually and ultimately is, as long as we are sure that the various forms of determination given us by progressing experimental and theoretical investigation can dovetail and mutually supplement one another}{So brauchen wir denn auf die Frage, was das Atom oder Elektron denn nun eigentlich und endgültig sei, keine Antwort zu finden, wenn wir nur sicher sind, daß die verschiedenen Formen seiner Bestimmung, die uns die fortschreitende experimentelle und theoretische Forschung in die Hand gibt, ineinander eingreifen und sich wechselseitig ergänzen} \DI{187}{150}. The \s{atom}, therefore, is not merely a subjective device for achieving economy of thought, as the positivists claim; but neither does it possess a thing-like existence, as the realist \cop{It is these historically stable properties that endure through subsequent theory changes. This accumulated stock of enduring properties can be collected into what we shall call the structure of electron theory}.\todo{cjeck}

\todo{stochimoetry and stereochmistry}. 


%That one recovers essentially the
%same value in all these circumstances is strong evidence that each is a mani-
%festation of the same particle, the electron, and that electrons do carry iner-
%tial mass and charge and in the ratio recovered. A second strategy applied to
%the electron is known in the philosophy of science literature by many names,
%including eliminative induction or demonstrative induction. In it, one maps
%out as large a class of candidate theories as possible and then shows that some
%item of evidence, usually experimental, forces selection of just one theory
%from that class as the only one that is compatible with this item of evidence.
%The force of this method is that it not only gives strong evidence for the the-
%ory selected, but it also gives direct evidence against the theory’s competitors.

%Indeed, modern quantum theory rejects the intuitive picture of the atom. In particular, the particle and wave pictures cannot be applied simultaneously, rguing that the list of its “historically stable” properties grew over time  preserving their successes and providing richer and improved representations of the electron namely, that over time the good drives out the bad, with what was true and correct in the past mov- ing forward, perhaps in a new form, into the future




%con ciò l’atomo ha tutta la « realtà » di cui è  capac


%Molecules with the same chemical formula can have different chemical properties. Atoms with the same atomic weight can have different chemical properties.

%They are distinct chemical compounds that share the exact same chemical formula They are atoms of the same element that have the exact same number of protons but a different number of neutrons



%Despite the philosophical disagreement, their personal relationship seemed to have remained excellent. In 1930, Meyerson asked Cassirer to sound out the possibility of German translation of the third revised edition of \IRt\footnote{Cassirer to Meyerson, Paris, 05.04.1930 (no. 621) enquires about the possibility of translating \IRt in German, Dienstag, den 8. April 1930 Signatur A408 / 8 \citep{Meyerson1930}}. At around the same time, Cassirer met Meyerson personally while visiting Paris for research. As his wife recounted, Cassirer \qt{was most warmly welcomed by the philosophers there---Lévy-Brühl, Brunschvicg, Meyerson, and others---and in fact developed much closer ties with this circle than with the German philosophers}{wurde von den dortigen Philosophen, Lévy-Brühl, Brunschvicg, Meyerson usw. aufs herzlichste aufgenommen und ist mit diesem Kreis eigentlich in viel engere Verbindung gekommen als mit den deutschen Philosophe} \citep[189]{CassirerBondy2003}. Cassirer, now on the verge of leaving Germany, did not seem to have taken into account Meyerson's later work \citep{Meyerson1931,Meyerson1933}. However, he clearly would not have seen their disagreement resolved in which he took stance on the new quantum theory. Meyerson once again attacked the widespread positivistic and phenomenalistic rhetoric of quantum physicists. Their attempt to dissolve matter into abstract mathematical equations was not a philosophical choice but a reluctant sacrifice (\textit{pis-aller, ersatz, renoncement}). In actual laboratory work and theoretical reasoning, the quantum physicists remain unwavering ontological realists. What Meyerson had called \german{chosisme} (\s{thingism}) reasserts itself as an inescapable law of human thought, as atoms and electrons are as real as tables and chairs. 

%\q{Man Does Metaphysics as He Breathes}.




%\qt{When asked wherein the strictly objective factor in nature consists, the modern physicist can only point to the universal constants arrived at by investigation and, on the other hand, trace the road leading from these universal constants down to the individual constants, the particular thing-constants}{Der moderne Physiker vermag auf die Frage, worin das eigentliche »Objektive« der Natur besteht, keine andere Antwort zu geben, als daß er die »universellen Konstanten« namhaft macht, bei deren Feststellung seine Forschung endet, und daß er andererseits den Weg von diesen universellen Konstanten bis zu den Individualkonstanten, zu den besonderen Dingkonstanten, hinab verfolgt} \PSF{}{}.


%objects of the second order

%Physicists tend to believe that an element's position in Mendeleev’s table—and its atomic number $n$—contains all the essential information needed to deduce its physical and chemical behaviors a priori [311, 760, 829–830].Under this view, chemical properties (valence, reactivity, spectral lines) cease to be irreducible qualitative "givens" and become spatial/quantitative consequences of electron configurations around the nucleus [829–830]. This reude owas reistatnce o reat, bt , that only true a mode. 


%Absolutum des naiven Dingbegriffs, sondern die Invarianz be- stimmter fundamentaler Gesetz Meyerson demonstrates that both everyday perception and advanced scientific theories operate through this exact same mental mechanism%he scientific intellect then constructs "second-order objects" (objets du deuxième ordre)—such as atoms, molecules, corpuscles, or electrons . These sub-sensory entities are projected into physical reality to serve as the persistent "things" that explain macroscopic transformation

%from the French word chose, meaning "thing"

%As one might expect, Meyerson draws the very opposite moral from this episode. In front a variey of chemical, elet  modern subatomic physics attempts to absorb chemistry and once to the intutieve model of Bohr's atom spatial mode \s{model} the number $n$ and spatial arrangement of electrons around a nucleus \citep[**]{Meyerson1931}. %The scientific intellect then constructs "second-order objects (\textit{objets du deuxième ordre}) such as atoms, molecules, corpuscles, or electrons,

%Yet, physicists like Niels Bohr and Louis de Broglie inevitably rely on symbolic analogies and visual particle images (billiard balls, orbiting bodies, wave models) to reason about subatomic phenomena. These sub-sensory entities are projected into physical reality to serve as the persistent "things" that explain macroscopic transformations. 

% Needless, to say Cassirer's once again draw opposite moral.


%Howevever, once again Even when modern physics attempts to strip away mechanical models and reduce the atom to abstract mathematical equations, configuration spaces, or wave functions, chosisme reasserts itself. 

%Cassirer did not seem to have however taken into account Meyerson's later work, neither nor his book on \qm from 1933 just before his death. Meyerson's \emph{choisisme} (as one might tranlate thingism) could not be foreign to Cassirer's philosophy. Form Myero Thre historu of chet conceot ofa tom 

%Cassirer remained in good with Meyerson to whom he to send some his last on relativity \letter{Cassirer}{Meyerson}{4}{7}{1921}[**][ECN]  



%Needless to say from the same episode Meyerson's draw the opposite conclusion, that a purely had ultimatle to retior intutieve model of Bohr's atom spatial mode \s{model} the number $n$ and spatial arrangement of electrons around a nucleus \citep[**]{Meyerson1931}. 


%Hoever, it soon becomes clear we we recong. Die Bilder versagen – aber die Relationen, die in den but numerical realation ramins deter. 

%Wir entdecken nunmehr, in den letzten empirischen »Naturgesetzen«, gleichsam Konstanten höherer Ordnung, die sich über dem bloß faktischen Bestand der Einzeltatsachen, der in bestimmten Größenwerten fixiert ist, erheben.



%Despite the diversity of accounts of the atom’s existence, nature, and structure, a direction exists through its unfolding sequence. It brings its particular forms together under a general rule, treating each as a distinct application. 


%We therefore need not establish what the atom or electron ultimately is, as long as the determinations emerging from experimental and theoretical research fit together and mutually complement one another. 


%Dann ist die Mannigfaltigkeit der Antworten, die wir auf die Frage nach dem »Sein« des Atoms, seiner Natur und Struktur erhalten, kein Einwand mehr. Denn eine Funktion ist nur dadurch, daß sie sich in eine fortschreitende Reihe von Werten entfaltet. Sie läßt sich nicht in ein einzelnes Bild pressen; sondern sie befaßt alle ihre Sonderformen unter einer universellen Regel, sie enthält sie als verschiedene Anwendungsfälle. Keiner dieser Anwendungsfälle ist ein endgültiges Ziel, an dem die Erkenntnis ausruhen könnte; er ist immer nur ein Meilenstein, an dem die Erkenntnis die Stelle, an der sie steht, und den Weg, den sie durchmessen, ablesen kann. So brauchen wir denn auf die Frage, was das Atom oder Elektron denn nun eigentlich und endgültig sei, keine Antwort zu finden, wenn wir nur sicher sind, daß die verschiedenen Formen seiner Bestimmung, die uns die fortschreitende experimentelle und theoretische Forschung in die Hand gibt, ineinander eingreifen und sich wechselseitig ergänzen

%Man sieht, wie es auch hier wieder der Zutritt neuer, sei es empirisch, sei es theoretisch geforderter Bedingungen ist, der dem Atom erst seine bestimmte »Gestalt« verleiht. Seine Struktur ist nichts anderes als der zusammenfassende Ausdruck dieser Bedingungen. An ihrer immer klareren und schärferen Formulierung sind alle großen Hauptgebiete der Physik beteiligt. Die Mechanik, die Chemie, die Thermodynamik, die Elektrodynamik greifen in gleicher Weise in diesen Prozeß ein – und indem sie sich in dieser Art auf ein gemeinsames Ziel richten, werden sie eben damit zu einer methodischen Einheit und Ganzheit zusammengeschlossen.



%For Meyerson the reduction of chemical atoms to physical atoms,  It is now sure that no fruitful attack on the problem could have been made before the discovery of the electron. W ebelieve in the 'existence' ofthe electron becauseit has adefinite charge e anda definite mass m and a definite spin s; that means in whatever circumstances andexperimentalconditions youobserveaneffect heory ascribes to the presence of electrons you find for these quantities, e, m, s, the same numerical values. First the attempt of field theory, that quantum mechanics the incvil of the 



%unktion des Atoms, als einer jeweilig letzten Maßeinheit, beharrt
%Here the physicists, going beyond chemistry, had laid hands on an elementary unit of matter which is the same in all chemically different substances and occurs freely in the form of cathode rays. The 'periodic system of elements' had previously pointed, though with a somewhat vague gesture, toward a uniform structure of the various chemical atoms. But now physics entered into its golden era of atomic research. Dur- ing the last half century it has provided a thorough and brilliant corro- boration for the basic tenets of atomism and penetrated into ever deeper layers of the strange atomic world. To begin with, all its methods led with increasing accuracy to the same values of the charge and mass of an electron. 



%Originally, elements were organized by how heavy they were. But weight is a continuous, analog measurement. Because elements exist as mixtures of different isotopes, atomic weight is actually just a statistical average The shift from atomic weight to atomic number was the profound realization that nature's ultimate organizing principle isn't based on how much something weighs, but on exactly how many discrete parts it has. 

%Indeed, the transformations of the concept of the atom provide a characteristic example of this continually advancing process. The atom, which in its genesis derives from the pure number-concept, has here, after manifold transformations, returned to its origin: it signifies nothing other than a member of a systematic sequence as such. Instead of studying each element in isolation, scientists realized that elemental properties are periodic functions of their atomic weight. Isotopes proved that atoms of the same element could have different atomic weights, meaning atomic weight was not a perfect, uniform metric for organizing the elements.  To resolve this complication, science required a new structural foundation. The text explains that atomic weight was merely a complex and unexplained function of a much simpler, truer metric, the \s{atomic number}: \q{Die »Nummer«, die ein Element im periodischen System getragen hatte, seine »Ordnungszahl« wird jetzt einfach durch die Zahl seiner Kernladung ersetzt – und diese gibt wiederum die Zahl der den Kern umgebenden Elektronen an (van den Broek 1913)}








%Eine neue Entwicklung des Atombegriffs setzt mit den Anfängen der wissenschaftlichen Chemie ein.  Die Forderung eines streng korpuskularen Aufbaus der Chemie wird bereits im 17. Jahrhundert durch Robert Boyle gestellt. Boyles »Chymista scepticus« (1677) führt die These durch, daß nur durch Einführung des Atombegriffs der Herrschaft der »dunklen Qualitäten« im Reich der Chemie ein Ende bereitet werden könne. \textelp{} Die Chemie verharrt zunächst noch innerhalb der qualitativen Naturbetrachtung, die in der Stahlschen Phlogistontheorie noch einmal in charakteristischer Schärfe hervortritt. Aber als dann durch Lavoisiers Versuche der Umschwung eintritt, als auch alle chemischen Tatsachen unter den Gesichtpunkt von Maß, Zahl und Gewicht gestellt werden – da ist auch hier die Stunde für die Atomistik gekommen. \textelp{} In Daltons Gesetz der multiplen Proportionen feiert die Atomistik ihren ersten Triumph; sie vermag unmittelbar die empirische Tatsache der konstanten Verbindungszahlen zu erklären. Die festen Äquivalentzahlen fanden ihre einfachste und sich unmittelbar aufdrängende Deutung in der Annahme, daß die Atome der verschiedenen Elemente ein bestimmtes Gewicht besitzen, das für die einzelnen Grundstoffe charakteristisch ist. Hierzu tritt | die Avogadrosche Regel, nach der die Anzahl der Moleküle in der Raumeinheit bei allen Gasen die gleiche ist. 






%The further conceptual determination advances, the less possible it becomes to capture its result in simple, individual intuitive images. We must content ourselves with determination through laws and forgo clarification through models and schemata. \q{Die wahre physikalische »Objektivität« aber ist dadurch nicht vermindert, sondern vermehrt; denn sie kann von Anfang an nicht in der Richtung auf den »naiven« Dingbegriff, sondern sie muß in der Richtung auf den kritischen Gesetzesbegriff hin gesucht werden.}. The concept of atom in chemistry:







%%The same process of progressive  By organizing elements by their atomic weight, chemistry found a "new heuristic principle."  This functional order transformed chemistry into a highly predictive science, allowing researchers to accurately forecast the existence and properties of undiscovered elements. Specifically radioactivity and the "isotope theory"—complicated this early model. 
%The Modern Paradigm (Atomic Number): T







% In dem Maße aber, als neben diese besondere Ausprägung des Atombegriffs nunmehr andere und andere treten und er schließlich die verschiedenartigsten sachlichen Gebiete beherrscht und gedanklich organisiert, tritt auch sein reiner Prinzipiencharakter, der anfangs mit seinem Sachcharakter noch unterschiedslos verschmolz, immer be|stimmter zutage.

%Die Geschichte des Begriffs der Materie, die Geschichte des Atombegriffs, des Begriffs des Äthers und der Energie liefert hierfür die typischen Beweise und Belege. Progressivel 

%More more and more material constants become connected through fixed laws, reducing the number that must be determined independently. Particular empirical relations are then incorporated into broader theories involving fewer, more universal constants. Initially, an element's atomic weight (a chemical property) and its specific heat (a thermodynamic property determining how much energy is needed to raise its temperature) were treated as separate constants. Dulong and Petit discovered they are intimately connected: they are inversely proportional. As measurements improved, a crisis emerged. The Dulong-Petit constant broke down for elements with low atomic weights (like carbon/diamond or boron). Their specific heats deviated significantly from what the kinetic theory predicted, opening a gap between theoretical expectations and experimental reality.  Einstein and Debye took the final step. By applying quantum mechanics, they proved that specific heat is not actually a static "material constant" at all. Instead, it is a variable that changes with absolute temperature, governed by universal, non-material constants (like Planck's constant and the Boltzmann constant). 



%Einstein resolved this anomaly by quantum of of action to the thermal vibrations of solid bodies, proving that the low-atomic-weight anomalies were mathematically predictable under quantum rules. Peter Debye refined this further, establishing a new, comprehensive law that linked specific heat directly to absolute temperature. \citep{Cassirer1929}.

%Here the physicists, going beyond chemistry, had laid hands on an elementary unit of matter which is the same in all chemically different substances and
%occurs freely in the form of cathode rays. The 'periodic system of
%elements' had previously pointed, though with a somewhat vague
%gesture, toward a uniform structure of the various chemical atoms.

%The old dream of the unity of all matter had certainly come a great step nearer to fulfillment by the discovery of the electron. But the positively charged nucleus of the atom about which the negatively charged electrons revolve like the planets around the sun still seemed to be a particle with a constitution of its own for each individual chemical element. It is a priori clear that beside the negative electron at least one positive brick is needed for the construction of all atoms. The execution of this idea of building up matter from two ultimate elementary units, the electron and the proton, presupposes of course that they are able to enter into a variety of combines that are held together by strong forces and react outwardly like solid atomic spheres. Hence sheer substance without force would never do.


%For it was possible to derive from
%observation in combination with the mechanical theory rather reliable
%values for the magnitude of the radii as well as for the inert masses of
%atoms, which enter into the expressions for energy and momentum;
%and it transpired that for the various chemical elements the atomic
%masses are far from proportional to the atomic volumes. This shat-
%tered the basic conception of one matter, the conception of a homo-
%geneous dough of substance out of which the Creator, with the help
%of a set of baking moulds, at the beginning of time had carved the
%little atom cakes, and had then given them absolute rigidity and sent
%them off into space with varying initial momenta. The mass ratios
%however, proved to be in accordance with the relative atomic weights,
%as derived from the actual quantitative analysis of innumerable chemi-
%cal compounds. The law of multiple proportions, on which the atom-
%istic interpretation of the findings of chemistry is based, was for a
%long time by far the most convincing empirical proof for the atomic
%constitution of matter.

%The field la'Ys should thus permit us to compute in
%advance charge and mass of the electron and the atomic weights of the
%various chemical elements in existence

%Die Atome werden von Demokrit als letzte Konstanten der Natur gefordert, lange bevor sich für das Denken irgendein Weg zur konkreten Erfüllung dieser Forderung zeigt. Im Grunde wird eine solche Erfüllung, wird ein streng exakter quantitativer Sinn des Atombegriffs erst mit den Anfängen der modernen Chemie im Gesetz der multiplen Proportionen erreicht.


%In fact, such a fulfillment—a rigorously exact quantitative meaning of the concept of the atom—is achieved only with the beginnings of modern chemistry, in the law of multiple proportions. Yet as this particular formulation of the concept of the atom is joined by ever further formulations, and as the concept ultimately comes to govern and conceptually organize the most diverse domains, its character as a pure principle, initially fused indistinguishably with its character as a thing, also emerges ever more clearly.

%\q{Im Grunde wird eine solche Erfüllung, wird ein streng exakter quantitativer Sinn des Atombegriffs erst mit den Anfängen der modernen Chemie im Gesetz der multiplen Proportionen erreicht. In dem Maße aber, als neben diese besondere Ausprägung des Atombegriffs nunmehr andere und andere treten und er schließlich die verschiedenartigsten sachlichen Gebiete beherrscht und gedanklich organisiert, tritt auch sein reiner Prinzipiencharakter, der anfangs mit seinem Sachcharakter noch unterschiedslos verschmolz, immer be|stimmter zutage} .



%
%
%
%
%%Einen wie selbständigen Anteil an dieser gesamten Entwicklung das letztere Moment besitzt, wie bedeutsam die Ausbildung einer wissenschaftlichen Formelsprache für die Aufstellung einer universellen Systematik der Naturgegenstände und der Naturvorgänge zu werden vermag, läßt sich auch von anderer Seite her erweisen. Die Chemie ist zur »exakten« Wissenschaft nicht nur durch die ständige Verfeinerung ihrer Maßmethoden geworden, sondern vor allem auch durch die Verschärfung dieses ihres gedanklichen Instruments, durch den Weg, den sie von der einfachen chemischen Formel bis zur Strukturformel durchmessen hat. Ganz allgemein besteht der wissenschaftliche Wert einer Formel nicht nur darin, daß sie gegebene empirische Tatbestände zusammenfaßt, sondern daß sie neue Tatbestände gewissermaßen hervorlockt. Sie stellt Probleme von Zusammenhängen, von Verknüpfungen und Reihenbildungen auf, die der unmittelbaren Beobachtung vorauseilen. So wird sie zu einem der hervorragendsten Mittel dessen, was Leibniz die »Logik der Entdeckung«, die logica inventionis, genannt hat. Schon die einfache chemische Formel, in der lediglich die Art der in einem bestimmten | 510	Molekül enthaltenen Atome sowie deren Anzahl angegeben wird, ist gleichsam erfüllt mit fruchtbaren systematischen Hinweisen. Wurden etwa, in der Sprache dieser Formel, gewisse bekannte Verbindungen von Chlor, von Wasserstoff und Sauerstoff als Cl O H, Cl O3 H, Cl O4H bezeichnet, so drängte sich schon in dieser bloßen Zusammenstellung die Frage nach dem »fehlenden Glied« dieser Reihe, nach einer Verbindung Cl O2 H auf, die nun erst, nachdem gewissermaßen ihr Ort im voraus bestimmt worden war, auch empirisch auffindbar wurde. Hier erweist sich der Erkenntniswert, der jeglicher methodisch aufgebauten wissenschaftlichen Sprache als solcher innewohnt. Sie ist niemals eine bloße Bezeichnung für Gegebenes und Vorhandenes, sondern ein Wegweiser in neues, bisher nicht erforschtes Gebiet; sie leitet zu einem Prozeß der »Interpolation« und »Extrapolation«. »Man ersieht | hieraus«, so schließt der Autor, dem ich das obige Beispiel entnehme, »wie sich in der chemischen Formelsprache, von Stufe zu Stufe, eine immer nähere Übereinstimmung zwischen den Symbolen und der Wirklichkeit herstellt und wie diese Sprache, die ursprünglich lediglich die Zusammensetzung der Stoffe gemäß ihren Gewichtsverhältnissen beschreiben sollte, nunmehr die Art ihrer Erschaffung beschreibt und uns den Einblick in sie aufzwingt. Es handelt sich nicht mehr um ein Verfahren der Bezeichnung, sondern um einen Leitfaden für die Entdeckung, um eine Methode der Synthese. […] Unsere Klassifikation gewinnt ein neues Gesicht: Sie bedeutet nicht länger eine Ordnung, die wir in der Art, wie sie uns die Natur oder der Zufall der Beobachtungen liefert, hinnehmen, sondern sie wird zu einer von uns selbst geschaffenen, zu einer deduktiven Ordnung. Und ebendies trägt dazu bei, der Chemie ihren eigentümlichen Charakter zu geben.«[193] Noch deutlicher tritt dieser Charakter zutage, wenn wir die chemische Formel auf derjenigen Stufe ihrer Entwicklung betrachten, auf der sie zur eigentlichen »Konstitutionsfor|mel«  511	wird. Eine Konstitutionsformel, wie diejenige, die Baeyer für das Indigo gegeben hat, setzt an Stelle der bloß empirischen Beschreibung einen echt genetischen Aufbau: Sie wird zu einer Aussage nicht nur über das »Daß«, sondern auch über das »Wie«, indem sie die Verbindung, um die es sich handelt, gleichsam vor unserm innern Auge erstehen läßt. Und daß diese Wendung zur »Genese« nicht nur ein vereinzeltes Motiv darstellt, das gleichsam sporadisch in besonderen Gebieten der Wissenschaft auftritt, sondern daß wir in ihr einen Grundzug der physikalischen und chemischen Begriffsbildung selbst zu sehen haben: dies wird unmittelbar deutlich, sobald wir die größten systematischen Errungenschaften betrachten, die dieser Begriffsbildung im Laufe des letzten Jahrhunderts beschieden gewesen sind. Die Aufstellung des »natürlichen Systems der Elemente«, wie sie im Jahre 1870 durch Lothar Meyer und Mendelejeff versucht wurde, bildet auch in rein gedanklicher und methodischer Hinsicht einen wichtigen Wendepunkt. Denn hier war schärfer und bewußter als zuvor die Forderung gestellt, die Mannigfaltigkeit der Elemente und die Verschiedenheit, die sich in ihren physikalischen und chemischen Eigenschaften ausdrückt, nicht einfach hinzunehmen, sondern einen Gesichtspunkt zu finden, von dem aus diese Mannigfaltigkeit sich übersehen und gliedern, sich nach einem festen | Prinzip ordnen ließ. Als solches ordnendes Prinzip wurde zunächst das Atomgewicht gewählt. Reiht man sämtliche bekannten Elemente nach wachsender Größe ihrer Verbindungsgewichte aneinander, so erhält innerhalb dieser Reihe jedes Element eine gewisse Stelle zugewiesen, die durch eine für dasselbe charakteristische Zahl, durch seine »Ordnungszahl«, bezeichnet wird. Bei dieser Gruppierung der Elemente nach ihren Ordnungszahlen zeigte sich, daß die wichtigsten Eigenschaften der Elemente eine eigentümliche Periodizität aufweisen; sie sind nicht wahllos über das Gesamtgebiet verstreut, sondern folgen in ihrem Auftreten einer festen Regel der Wiederkehr. Die bekannte Kurve der Atomvolumina, wie sie Lothar Meyer aufgestellt hat, zeigt in anschaulicher Deutlichkeit das Bestehen und die Art dieser Regel. Die Elemente, die auf analogen Stellen dieser Kurve – auf ihrem aufsteigenden oder absteigenden Ast, bei einem ihrer Maxima oder Minima – liegen, erweisen sich auch in ihren wichtigsten chemischen und physikalischen Eigenschaften, in ihrer Valenz, ihrer Flüchtigkeit, Dehnbarkeit, elektrischen und thermischen Leitungsfähigkeit usf. als einander analog. Daß diese Abhängigkeit der Eigenschaften von der »Ordnungszahl« einen tieferen systematischen »Grund« haben, daß sie irgendwie in der »Natur« des Atoms selbst verankert sein müsse, | 512	ist eine Voraussetzung, von der schon die ersten Entdecker des periodischen Systems beherrscht waren und die ihnen als das eigentliche Motiv, als die heuristische Maxime ihrer Forschung diente. Aber diese Maxime war anfangs noch weit entfernt von einer eigentlichen »konstitutiven« Bedeutung. Denn der Zusammenhang zwischen dem Atomgewicht oder Atomvolumen eines Elements und seinem besonderen chemisch-physikalischen Verhalten war selbst zunächst lediglich als Tatsache festgestellt, ohne daß sich von dieser Tatsache eine wahrhaft befriedigende theoretische Rechenschaft geben ließ. So fungierten in den tabellarischen Darstellungen des natürlichen Systems die Ordnungszahlen, die jedem Element beigegeben wurden, zunächst einfach als konventionelle Marken, von denen sich innerhalb bestimmter Grenzen die Eigenschaften der Elemente ablesen ließen, die aber noch keinen bestimmten physikalischen Sinn in sich schlossen. Der weitere Fortgang der Theorie aber bestand nun eben darin, daß in immer höherem Maße dieser »Sinn« erkannt und herausgearbeitet und daß eben damit erst die konventionelle Ordnung in eine eigentlich systematische umgeprägt wurde. Der erste Schritt auf diesem Wege bestand in der Gewinnung eines neuen, schärferen Ordnungsprinzips, das sich der Spektroskopie der Röntgenstrahlen entnehmen ließ. Versuchte man | die verschiedenen Elemente nach ihren charakteristischen Röntgenspektren in eine Reihe zu ordnen, so ergab sich, daß innerhalb dieser Reihe die Verschiebung der einzelnen Linien von einem Element zum andern in der Richtung zunehmender Schwingungszahl mit einer Gesetzmäßigkeit erfolgt, die der Regelmäßigkeit, die sich in der Aufreihung nach zunehmenden Atomgewichten ergeben hatte, bei weitem überlegen war. Nach dem von Moseley im Jahre 1913 aufgestellten Gesetz ändert sich die Quadratzahl aus der Schwingungszahl einer charakteristischen Röntgenlinie fast genau linear mit der Ordnungszahl des chemischen Elements. Diese Tatsache ließ von Anfang an eine tiefe physikalische Bedeutung ebendieser Ordnungszahl vermuten: Denn nach den modernen theoretischen Grundanschauungen ist der Ort für die Entstehung des Röntgenspektrums das eigentliche Atominnere, der Kern des Atoms, während die optischen Spektren sowie die chemischen Eigentümlichkeiten auf mehr »äußerlichen«, auf peripheren Eigenschaften des Atoms beruhen.[194] Die Bekanntschaft mit weiteren empi|rischen  513	Tatsachen, insbesondere die Entdeckung der sogenannten »Isotopen« und die Ausbildung der Isotopentheorie, führten sodann dazu, daß das Atomgewicht, von dem die Gliederung des natürlichen Systems zunächst ausgegangen war, seine beherrschende Stellung in der Chemie verlor und daß diese nunmehr an einen anderen Begriff, an den Begriff der »Kernladung«, überging. Schon Moseley selbst hatte diesen Begriff in den Mittelpunkt der Systematik gestellt: Er hatte die »Ordnungszahl« einfach als diejenige Zahl gedeutet, die die positive Ladung des Atomkerns angibt. Diese Größe der Kernladung erscheint nunmehr als das eigentliche und endgültige Reihenprinzip. Die »Nummer«, die ein Element im periodischen System getragen hatte, seine »Ordnungszahl« wird jetzt einfach durch die Zahl seiner Kernladung ersetzt – und diese gibt wiederum die Zahl der den Kern umgebenden Elektronen an (van den Broek 1913). Als die gesuchte physikalische Bedeutung der »Ordnungszahl« ergibt sich somit die Annahme, die die Grundlage der Bohrschen Atomtheorie bildet, daß die positive elektrische Ladung des Atomkerns von Element zu Element je um eine Einheit zunimmt. »Der Atombau wird auf elektrischem Wege einheitlich von innen heraus bis an die Atomperipherie hin durch die | Größe der Kernladung geregelt.«[195] »Wir haben ein Ziel erreicht«, so faßt Sommerfeld in seinem Werk über »Atombau und Spektrallinien« den geistigen Ertrag dieser Entwicklung zusammen, »das vor zehn Jahren noch in nebelhafter Ferne zu liegen schien: eine Theorie des periodischen Systems.«[196] Betrachtet man den rein logischen Charakter dieser Theorie, so erkennt man, wie die Leistung, die in ihr vollbracht ist, sich zwar streng im Kreise des Beobachtbaren, des empirischen Feststellbaren hält – wie sie aber nichtsdestoweniger über jenen resignierten »Empirismus«, wie ihn Locke in seiner Substanzenlehre vertreten hatte,[197] weit hinausgeht. Sie definiert die Substanz keineswegs als ein bloßes Aggregat von Eigenschaften, die durch kein »inneres Band« miteinander verknüpft sind – aber sie faßt andererseits dieses innere Band nicht als ein »vinculum substantiale« im Sinne einer dogmatischen Metaphysik, sondern sie fragt lediglich nach der »Notwendigkeit« der Verknüpfung im Sinne ihrer durchgreifenden Allgemeinheit und Gesetzlichkeit. Diese Ge|setzlichkeit  514	ist niemals durch die Beobachtung vereinzelter »Tatsachen« zu gewinnen, so weit diese auch getrieben werden mag, sondern ihre Feststellung bedarf jederzeit bestimmter konstruktiver Gesichtspunkte und konstruktiver Prinzipien. Sowenig diese Gesichtspunkte dem empirischen Material aufgezwungen werden können, sowenig können sie ihm selber unmittelbar entnommen werden. Die Begriffsbildung der Physik und Chemie erweist sich daher ebenso als echt »genetische« Begriffsbildung, wie dies innerhalb der reinen Mathematik galt. Aber die Genesis, auf die hier ausgegangen wird, ist gewissermaßen nicht kategorischer, sondern hypothetischer Art. Wir beginnen nicht mit dem allgemeinen Reihengesetz, um aus ihm die Mannigfaltigkeit der Elemente hervorgehen zu lassen, sondern wir begnügen uns damit, der »gegebenen« Mannigfaltigkeit in verschiedenen gedanklichen Ansätzen versuchsweise ein ordnendes Prinzip zu unterbreiten, um damit Schritt für Schritt die bloß empirische Vielheit in eine »rationale« umzugestalten.[198] | Dieses Prinzip selber ist uns niemals schlechthin »gegeben«, wohl aber ist es beständig »aufgegeben« – und in der immer vollkommeneren Lösung dieser Aufgabe besteht eine der wesentlichen Leistungen aller Naturtheorie.
%
%
%%Perrin’s strongest argument was the agreement of different experimental methods on the same molecular quantity: Avogadro’s number. Brownian motion was central because it made the effects of otherwise invisible molecules measurable. 
%%
%%Perrin’s strongest argument was the agreement of different experimental methods on the same molecular quantity: Avogadro’s number. Brownian motion was central because it made the effects of otherwise invisible molecules measurable. For Cassirer, different elements exemplify the same quantitative relation between properties. The distinction is therefore better expressed as agreement among independent determinations versus invariance across different cases. A Cassirerian reading of Perrin would locate the empirical warrant for atoms in the convergence of independent experimental determinations. What this convergence establishes is the atom’s place within a system of measurable relations; interpreting it as a substantial entity existing behind those relations constitutes a further commitment.
%
%% Immer handelt es sich für die physikalische Auffassung darum, zunächst einen charakteristischen Gesichtspunkt der Messung in einem objektiv-physikalischen Begriff, in einer bestimmten Naturkonstante, zu fixieren – dann aber im weiteren Verlauf das konstruktive Moment, das in jeder solchen Setzung einer ursprünglichen Konstante liegt, klarer und klarer zu durchschauen und sich seiner Bedingtheit bewußt zu werden. Denn alle Konstanten, wie immer sie im einzelnen beschaffen sein mögen, werden nicht unmittelbar gegeben, sondern sie müssen zuvor gedacht und gesucht sein, ehe sie in der Erfahrung gefunden werden können. Eines der prägnantesten Beispiele hierfür steht in der Geschichte des Atombegriffs vor uns. Die Atome werden von Demokrit als letzte Konstanten der Natur gefordert, lange bevor sich für das Denken irgendein Weg zur konkreten Erfüllung dieser Forderung zeigt. Im Grunde wird eine solche Erfüllung, wird ein streng exakter quantitativer Sinn des Atombegriffs erst mit den Anfängen der modernen Chemie im Gesetz der multiplen Proportionen erreicht. In dem Maße aber, als neben diese besondere Ausprägung des Atombegriffs nunmehr andere und andere treten und er schließlich die verschiedenartigsten sachlichen Gebiete beherrscht und gedanklich organisiert, tritt auch sein reiner Prinzipiencharakter, der anfangs mit seinem Sachcharakter noch unterschiedslos verschmolz, immer be|stimmter zutage. Der Inhalt der Atomvorstellung wandelt sich und rückt im Lauf der Entwicklung der Physik und Chemie von Ort zu Ort – aber die Funktion des Atoms, als einer jeweilig letzten Maßeinheit, beharrt. Wenn wir von der Betrachtung der »ponderablen« Materie zur Betrachtung des Äthers übergehen, wenn wir eine Einheit suchen, die nicht nur die mechanischen, sondern auch die optischen und elektromagnetischen Phänomene umfaßt – dann gestaltet sich für uns das Atom der Materie zum Atom der Elektrizität, zum Elektron. In der neueren Physik tritt ferner mit der Planckschen Quantentheorie der Gedanke einer atomistischen Struktur nicht nur der Materie, sondern auch der Energie hervor. Es wäre ein vergebliches Bemühen, wollte man versuchen, alle diese verschiedenen Anwendungen, die der Atomgedanke in der Chemie, in der kinetischen Gastheorie, in der Lehre von der Licht- und Wärmestrahlung usf. erfahren hat, zu einem einheitlichen Bilde zusammenzufassen. Aber die Einheit seiner Bedeutung bedarf und fordert auch keine solche Bildeinheit; sie ist zur Genüge, ja in einem weit strengeren logischen Sinne gewahrt, wenn sich herausstellt, daß hier überall eine gemeinsame Beziehung, eine eigentümliche »Form« des Zusammenhangs, obwaltet, die sich als solche an den ver|schiedenartigsten  12	Inhalten bewähren und darstellen kann. Das Atom erweist sich eben darin nicht als ein absolutes Minimum des Seins, sondern als ein relatives Minimum des Maßes. Es ist einer der Begründer der neueren Philosophie, Nikolaus Cusanus, gewesen, der diese seine Aufgabe und Funktion, die es erst in der Geschichte der Naturwissenschaft wahrhaft erfüllen sollte, mit echtem spekulativem Tiefsinn vorausgesehen und verkündet hat. Cusas Grundlehre vom Unendlichen und von der Einheit der Gegensätze im Unendlichen ruht ganz auf dieser Einsicht in die prinzipielle Relativität aller Größenbestimmung – auf der Koinzidenz des »Größten« und »Kleinsten«.[12] Die moderne Erkenntniskritik kann die Rätsel, mit denen Cusas Lehre vom Minimum ringt, in einem einfachen Grundgedanken zusammenfassen. Der Widerspruch tritt nur ein, wenn wir alle die verschiedenen Formen, die der Gedanke des »Kleinsten«, der Gedanke der letzten unteilbaren Maßeinheit, in den verschiedenen Gebieten der Betrachtung annimmt, dinglich in eins zu fassen versuchen; aber er verschwindet sogleich, wenn wir uns darauf besinnen, daß die wahre Einheit niemals in den Dingen als solchen, sondern in den gedanklichen Setzungen | zu suchen ist, die wir je nach der Besonderheit des Gebiets, das es auszumessen und zu bestimmen gilt, wählen und die daher prinzipiell einer unbegrenzten Variabilität fähig sind. Es ergibt sich hieraus, daß, wie das Gemessene an Zahl unbegrenzt ist, so auch das Messende in unendlich vielen und unendlich verschiedenen Ansätzen bestehen und sich darstellen kann – die Einheit, die wir zu suchen und zu fordern haben, liegt weder in dem einen noch in dem anderen Glied, sondern lediglich in der Form ihrer wechselseitigen Verknüpfung, d. h. in den logischen Bedingungen der Operation der Messung selbst.
%
%
%%Als Beispiel hierfür kann die Geschichte des chemischen Elementenbegriffs dienen, in der dieser Prozeß sich in höchst bezeichnender und bedeutsamer Weise darstellt. In den Anfängen der wissenschaftlichen Chemie tritt der Begriff des Elements als ein Fundamentalbegriff auf, der für den Aufbau ihres theoretischen Systems unentbehrlich und von entscheidender Bedeutung ist. Aber der Umfang des Begriffs scheint sich zunächst keinerlei festen Regel unterwerfen zu lassen; sowohl die Zahl der Elemente wie ihre Anordnung bleibt schlechthin »zufällig«. Bis zum Schluß des 19. Jahrhunderts waren etwa 70 chemische Elemente bekannt, deren Zahl sich dann schnell erweitert hat; sie wird gegenwärtig als 92 angenommen. Was die Klassifikation dieser Elemente betraf, so mußte sich die Wissenschaft zunächst damit begnügen, sie unter gewissen, relativ äußerlichen Gesichtspunkten zu versuchen: Man teilte etwa die Elemente in Metalle und Metalloide ein. Eine scharfe Grenze ließ sich aber hierbei nicht ziehen: Es gab Grundstoffe, die ihrem äußeren Ansehen nach als Metalle erschienen, sich aber in chemischer Beziehung durchaus wie Metalloide verhielten – und umgekehrt. Ein wesentlicher Fortschritt und ein Umschwung der gesamten Betrachtungsweise wurde erst erreicht, als durch Lothar Meyer und Mendelejeff, im Jahre 1871, das periodische System der Elemente aufgestellt wurde, das auf einer Gruppierung der Elemente | 26	nach der Größe ihrer Atomgewichte beruht. Hier war mit einem Schlage ein neues heuristisches Prinzip gefunden, auf das die Klassifikation sich stützen und kraft dessen sie eine unvergleichlich größere Strenge und Sicherheit erreichen konnte. Die Eigenschaften der Elemente und ihrer Verbindungen brauchten nicht länger einzeln, von Fall zu Fall, festgestellt zu werden; sie ergaben sich in gewisser Weise aus dem Inhalt des Elementenbegriffs, wenn man in diesen Inhalt das Moment des Atomgewichts aufnahm. Denn jetzt war eine funktionale Ordnung gefunden: Die Eigenschaften erschienen als periodische Funktionen der Atomgewichte. Ordnete man die Elemente nach steigendem Atomgewicht, so fand sich, daß nach gewissen Intervallen ähnliche Elemente wiederkehrten. Nicht nur ließen sich jetzt fast alle physikalischen und chemischen Eigenschaften eines gewissen Elements durch seine Stellung im System bestimmen – es zeigte sich auch, daß es auf Grund dieses Systems möglich war, die Existenz bestimmter noch nicht beobachteter Grundstoffe vorauszusehen und diese Voraussage durch die Erfahrung zu bestätigen. Ein weiterer, empirisch und methodisch gleich bedeutsamer Fortschritt vollzog sich sodann auf Grund der neuen Daten, die das Studium der radioaktiven Erscheinungen lieferte. Durch die Erkenntnis der radioaktiven Verschiebungs|sätze und durch die »Isotopentheorie« wurde man zu der Einsicht geführt, daß das periodische System nicht von so einfacher Periodizität war, als man bisher angenommen hatte. Es mußte ein neuer Aufbau gesucht werden, der auf einem viel umfassenderen und komplizierteren Tatbestand beruhte, der aber andererseits, eben durch die Berücksichtigung dieses Tatbestandes, zu einem überraschend einfachen Ordnungsprinzip hinführte. Das Atomgewicht konnte nicht länger als das wahre ordnende Prinzip im natürlichen System festgehalten werden; es erwies sich, daß es nur eine verwickelte und noch ungeklärte Funktion der wahren »Ordnungszahl« darstellte. Als diese wahre Ordnungszahl wurde die Anzahl der positiven Ladungseinheiten des Kerns oder die Elektronenzahl des Atoms gefunden.[32] Damit war in der Tat alles erreicht, was innerhalb einer empirischen Theorie erreichbar ist und was eine solche sich als Ziel ihrer Begriffsbildung stecken kann. Daß jetzt, innerhalb des Umfangs des chemischen Elementenbegriffs, von einer »gleichgültigen Zufälligkeit«, einer »unbestimmbaren Regellosigkeit« nicht mehr die Rede sein konnte, liegt auf der Hand. Aber auf der | 27	anderen Seite zeigte sich, daß die Regel in keiner Weise antizipiert, daß sie nicht aus dem einfachen spekulativen »Begriff der Sache« entnommen werden konnte. Erfahrung und Denken mußten sich hier das Gleichgewicht halten und sich ständig wechselseitig ergänzen. Die Theorie gestaltet, an der Hand der Beobachtung und unter ihrer stetigen Leitung, den Inhalt der Begriffe um, indem sie immer neue Bestimmungen und Beziehungen in sie aufnimmt; und eben durch diese Umgestaltung wird sie fähig, die einzelnen Anwendungsfälle der Begriffe immer vollständiger zu übersehen und immer genauer zu ordnen. Der Umfang hört durch diese Ordnung auf, ein bloßes Aggregat zu sein – er wird zum System; zu einem Ganzen, das nicht nur durch den Begriff umfaßt, sondern das von einem bestimmten Mittelpunkt aus verstanden und gegliedert wird. Das anfängliche lockere »Nebeneinander« der Elemente wird zu einem »Ineinander«, das gleichwohl nur als ein empirischer Zusammenhang angesehen werden darf und daher immer nur relativ zur jeweilig erreichten Stufe der Erfahrungserkenntnis seine Gültigkeit hat. 
%
%
%% Cassirer, however, emphasizes a different kind of convergence: distinct empirical laws jointly constrain the determination of atomic weights. Three are especially important in his account: Avogadro’s rule, relating gas volumes and vapor densities to molecular weights; the Dulong--Petit law, relating atomic weights to specific heats; and Mitscherlich’s law of isomorphism, relating crystalline form to chemical composition. Here, the evidential force lies in the mutual agreement of determinations drawn from different domains of physical and chemical phenomena.
%
%%For Perrin, different experimental methods converge on the value of the same molecular constant.
%
%
%
%%Eine neue Entwicklung des Atombegriffs setzt mit den Anfängen der wissenschaftlichen Chemie ein. Die Forderung eines streng korpuskularen Aufbaus der Chemie wird bereits im 17. Jahrhundert durch Robert Boyle gestellt. Boyles »Chymista scepticus« (1677) führt die These durch, daß nur durch Einführung des Atombegriffs der Herrschaft der »dunklen Qualitäten« im Reich der Chemie ein Ende bereitet werden könne. Der Schritt von der Alchimie zur Chemie ist in diesem Werke prinzipiell vollzogen; es lehrt, daß wir die chemischen Tatsachen nicht durch die Annahme geheimer »Sympathien« zwischen den Körpern, sondern nur aus den Struktur- und Bewegungsgesetzen ihrer kleinsten Teile verstehen können. Aber noch bleibt diese Anschauung unfruchtbar, da die nötigen quantitativen Ansätze zu ihrer Durchführung fehlen. Die Chemie verharrt zunächst noch innerhalb der qualitativen Naturbetrachtung, die in der Stahlschen Phlogistontheorie noch einmal in charakteristischer Schärfe hervortritt. Aber als dann durch Lavoisiers Versuche der Umschwung eintritt, als auch alle chemischen Tatsachen unter den Gesichtpunkt von Maß, Zahl und Gewicht gestellt werden – da ist auch hier die Stunde für die Atomistik gekommen. Lavoisier erklärt, daß alle chemische Experimentierkunst auf dem Grundsatz beruhe, |
%178	daß eine wahrhaftige Gleichheit oder Gleichung zwischen den anfänglich gegebenen Körpern und denen, die man aus ihnen durch die Analyse gewinnen kann, besteht.[73] So müssen auch alle Aussagen über die Atome, zu denen diese Experimentierkunst führt, zuletzt dazu dienen, dieser »Gleichheit« einen bestimmten Ausdruck zu geben und sie theoretisch verständlich zu machen. In Daltons Gesetz der multiplen Proportionen feiert die Atomistik ihren ersten Triumph; sie vermag unmittelbar die empirische Tatsache der konstanten Verbindungszahlen zu erklären. Die festen Äquivalentzahlen fanden ihre einfachste und sich unmittelbar aufdrängende Deutung in der Annahme, daß die Atome der verschiedenen Elemente ein bestimmtes Gewicht besitzen, das für die einzelnen Grundstoffe charakteristisch ist. Hierzu tritt | die Avogadrosche Regel, nach der die Anzahl der Moleküle in der Raumeinheit bei allen Gasen die gleiche ist. Auf dem einfachen Grundschema, das damit gewonnen ist, baut sich die gesamte folgende Entwicklung der Chemie auf.[74] Wie bei Huyghens, so wird hier das Atom streng substantiell gefaßt; aber alle Aussagen über sein »Was« gehen zuletzt in Aussagen über ein »Wieviel«, in Aussagen über bestimmte numerische Verhältnisse über. Alle Grundbegriffe und alle Theorien der wissenschaftlichen Chemie – der Valenzbegriff und der Typenbegriff wie der Begriff des Radikals und die Theorie der zusammengesetzten Radikale – weisen in diese Richtung. Die Begriffssprache der Atomistik erlaubt jetzt, in Gemeinschaft mit der allgemeinen Zeichensprache der Mathematik, jenen einfachen und großartigen Aufbau der chemischen Formelsprache, wie er sich im Fortgang von den einfachen Formeln zu den Strukturformeln vollzieht. Die Stereochemie, wie sie von van’t Hoff begründet wird, erlaubt auf Grund einer komplexeren Gestalt der Formeln eine neue Gruppe von Tatsachen, insbesondere die Erscheinungen der Isomerie, zum Ausdruck zu bringen; sie gibt einen neuen Aufschluß über die Lagerung der Atome im Raume. So wandelt und verfeinert sich das anschauliche Bild des Atoms; aber alle diese Wandlungen dienen nur dem einen Zweck, einen immer größeren Kreis funktioneller Beziehungen zu umspannen und darzustellen. 
%
%Each element has a characteristic assigned atomic weight in the historical framework Cassirer discusses. The amounts of that element entering different molecules occur in integer multiples of its atomic weight. These relations leave the molecular scale undetermined (p. 268). They do not establish \(N\), and therefore do not determine individual atomic masses in grams. Independent physical methods then determine that scale (pp. 277–281). Brownian motion, electrical measurements, radioactivity, and other phenomena yield mutually compatible determinations of \(N\). His concluding emphasis is precisely on phenomena that \s{convergent ainsi précisément vers le même nombre} (p. 281).  Perrin adds a determination that relative atomic weights alone cannot supply—the individual molecular mass scale. He expressly takes this as evidence for the actual existence of atom.
%
%
%
%The advance available to Cassirer in 1910 was from relative atomic weights to the determination of individual atomic and molecular masses.
%Relative weights establish that an oxygen atom is approximately sixteen times as massive as a hydrogen atom. Determining \(N_A\) additionally establishes their individual masses in grams, has led to establishing the
%\s{absolute} size of atoms. However, it seems that Cassirer that this additional determination enriches the system of quantitative relations through which the atom is objectively specified. Even this individual mass is a quantitative determination—a further Wieviel, a furhter numeical ccaratectit. That  Whether Cassirer explicitly understood Perrin’s contribution in precisely these terms would require evidence from his discussion or correspondence \q{so wird hier das Atom streng substantiell gefaßt; aber alle Aussagen über sein »Was« gehen zuletzt in Aussagen über ein »Wieviel«, in Aussagen über bestimmte numerische Verhältnisse über}.  Alle Grundbegriffe und alle Theorien der wissenschaftlichen Chemie – der Valenzbegriff und der Typenbegriff wie der Begriff des Radikals und die Theorie der zusammengesetzten Radikale – weisen in diese Richtung
%
%
%
%%Er erklärt, daß durch die Isolierung und Zählung der Gasionen, wie sie Joseph John Thomson gelungen sei, durch die Übereinstimmung der Brownschen Bewegung mit den Forderungen der kinetischen Hypothese usf. der experimentelle Beweis für die atomistische Beschaffenheit der raumerfüllenden Stoffe erbracht sei; aus einer bloßen Hypothese sei die Atomistik nunmehr zum Range einer wissenschaft|lich wohlbegründeten Theorie aufgestiegen
%
%
%
%
%
%
%%Die »Natur« eines Körpers, im physikalischen Sinne des Wortes, wird nicht durch die Art seiner sinnlichen Erscheinung, sondern durch sein Atomgewicht, seine spezifische Wärme, seinen Brechungsexponenten, seinen Absorptionsindex, seine elektrische Leitungsfähigkeit, seine magnetische Suszeptibilität usf. bestimmt
%
%
%Das Atom erweist sich eben darin nicht als ein absolutes Minimum des Seins, sondern als ein relatives Minimum des Maßes.
%
%%Im Grunde wird eine solche Erfüllung, wird ein streng exakter quantitativer Sinn des Atombegriffs erst mit den Anfängen der modernen Chemie im Gesetz der multiplen Proportionen erreicht. In dem Maße aber, als neben diese besondere Ausprägung des Atombegriffs nunmehr andere und andere treten und er schließlich die verschiedenartigsten sachlichen Gebiete beherrscht und gedanklich organisiert, tritt auch sein reiner Prinzipiencharakter, der anfangs mit seinem Sachcharakter noch unterschiedslos verschmolz, immer be|stimmter zutage. Der Inhalt der Atomvorstellung wandelt sich und rückt im Lauf der Entwicklung der Physik und Chemie von Ort zu Ort – aber die Funktion des Atoms, als einer jeweilig letzten Maßeinheit, beharrt
%
%Indem die Zahl dieser Bedingungsgleichungen wächst, nimmt das Atom in gewissem Sinne eine immer »konkretere« Gestalt an, weil eine immer größere Zahl besonderer Erkenntnisse sich in ihm vereinigt. Aber dieses Zusammenwachsen, diese »Konkreszenz« der Bestimmungsstücke steht zu dem, was wir unter anschaulicher Konkretion zu verstehen pflegen, im umgekehrten Verhältnis. Je weiter die |
%174	begriffliche Bestimmung fortschreitet, um so weniger erweist es sich als möglich, ihr Ergebnis noch in einfachen und einzelnen anschaulichen Bildern festzuhalten.
%
%Die wahre physikalische »Objektivität« aber ist dadurch nicht vermindert, sondern vermehrt; denn sie kann von Anfang an nicht in der Richtung auf den »naiven« Dingbegriff, sondern sie muß in der Richtung auf den kritischen Gesetzesbegriff hin gesucht werden. 
%
%errin's
%absolute magnitUdes of atom
%
%
%
%
%
%
%
%
%%Eines der prägnantesten Beispiele hierfür steht in der Geschichte des Atombegriffs vor uns. Die Atome werden von Demokrit als letzte Konstanten der Natur gefordert, lange bevor sich für das Denken irgendein Weg zur konkreten Erfüllung dieser Forderung zeigt. Im Grunde wird eine solche Erfüllung, wird ein streng exakter quantitativer Sinn des Atombegriffs erst mit den Anfängen der modernen Chemie im Gesetz der multiplen Proportionen erreicht. In dem Maße aber, als neben diese besondere Ausprägung des Atombegriffs nunmehr andere und andere treten und er schließlich die verschiedenartigsten sachlichen Gebiete beherrscht und gedanklich organisiert, tritt auch sein reiner Prinzipiencharakter, der anfangs mit seinem Sachcharakter noch unterschiedslos verschmolz, immer be|stimmter zutage. Der Inhalt der Atomvorstellung wandelt sich und rückt im Lauf der Entwicklung der Physik und Chemie von Ort zu Ort – aber die Funktion des Atoms, als einer jeweilig letzten Maßeinheit, beharrt. Wenn wir von der Betrachtung der »ponderablen« Materie zur Betrachtung des Äthers übergehen, wenn wir eine Einheit suchen, die nicht nur die mechanischen, sondern auch die optischen und elektromagnetischen Phänomene umfaßt – dann gestaltet sich für uns das Atom der Materie zum Atom der Elektrizität, zum Elektron. In der neueren Physik tritt ferner mit der Planckschen Quantentheorie der Gedanke einer atomistischen Struktur nicht nur der Materie, sondern auch der Energie hervor. Es wäre ein vergebliches Bemühen, wollte man versuchen, alle diese verschiedenen Anwendungen, die der Atomgedanke in der Chemie, in der kinetischen Gastheorie, in der Lehre von der Licht- und Wärmestrahlung usf. erfahren hat, zu einem einheitlichen Bilde zusammenzufassen. Aber die Einheit seiner Bedeutung bedarf und fordert auch keine solche Bildeinheit; sie ist zur Genüge, ja in einem weit strengeren logischen Sinne gewahrt, wenn sich herausstellt, daß hier überall eine gemeinsame Beziehung, eine eigentümliche »Form« des Zusammenhangs, obwaltet, die sich als solche an den ver|schiedenartigsten 
%%12	Inhalten bewähren und darstellen kann.
%
%
%%edes neue Gebiet, in dem die Atomistik zur Geltung kommt, liefert einen neuen Inbegriff von Bedingungsgleichungen, denen fortan das Atom gehorchen muß. Indem die Zahl dieser Bedingungsgleichungen wächst, nimmt das Atom in gewissem Sinne eine immer »konkretere« Gestalt an, weil eine immer größere Zahl besonderer Erkenntnisse sich in ihm vereinigt. Aber dieses Zusammenwachsen, diese »Konkreszenz« der Bestimmungsstücke steht zu dem, was wir unter anschaulicher Konkretion zu verstehen pflegen, im umgekehrten Verhältnis. Je weiter die | 174	begriffliche Bestimmung fortschreitet, um so weniger erweist es sich als möglich, ihr Ergebnis noch in einfachen und einzelnen anschaulichen Bildern festzuhalten. Wir müssen uns mit der Determination durch Gesetze genügen lassen und auf die Verdeutlichung durch Modelle und Schemata verzichten. Die wahre physikalische »Objektivität« aber ist dadurch nicht vermindert, sondern vermehrt; denn sie kann von Anfang an nicht in der Richtung auf den »naiven« Dingbegriff, sondern sie muß in der Richtung auf den kritischen Gesetzesbegriff hin gesucht werden
%
%
%
%%ines der wichtigsten Kennzeichen bietet hier die Avogadrosche Regel, nach welcher gleiche molekulare Massen verschiedener Verbindungen im vollkommenen Gaszustand unter denselben Bedingungen des Drucks und der Temperatur das gleiche Volumen einnehmen. Neben die Bestimmung der Atomgewichte aus der Dampfdichte, die hierdurch ermöglicht wird, tritt sodann ihre Bestimmung aus der Wärmekapazität, die sich auf das Dulong-Petitsche Gesetz, sowie diejenige aus dem Isomorphismus, die sich auf den Mitscherlichschen Satz stützt, daß die gleiche Kristallform, die verschiedene Verbindungen besitzt, auf eine gleiche Anzahl auf gleiche Art verbundener Atome in ihnen zurückweist. Erst die Gesamtheit aller dieser verschiedenen Gesichtspunkte, die sich wechselseitig bestätigen oder korrigieren, ergibt schließlich nach mannigfachen Versuchen eine einheitliche Tabelle der Atomgewichte und legt damit den Grund zu einer eindeutigen chemischen Formelsprache
%
%
%
%
%
%
%%Eine neue Entwicklung des Atombegriffs setzt mit den Anfängen der wissenschaftlichen Chemie ein. Die Forderung eines streng korpuskularen Aufbaus der Chemie wird bereits im 17. Jahrhundert durch Robert Boyle gestellt. Boyles »Chymista scepticus« (1677) führt die These durch, daß nur durch Einführung des Atombegriffs der Herrschaft der »dunklen Qualitäten« im Reich der Chemie ein Ende bereitet werden könne. Der Schritt von der Alchimie zur Chemie ist in diesem Werke prinzipiell vollzogen; es lehrt, daß wir die chemischen Tatsachen nicht durch die Annahme geheimer »Sympathien« zwischen den Körpern, sondern nur aus den Struktur- und Bewegungsgesetzen ihrer kleinsten Teile verstehen können. Aber noch bleibt diese Anschauung unfruchtbar, da die nötigen quantitativen Ansätze zu ihrer Durchführung fehlen. Die Chemie verharrt zunächst noch innerhalb der qualitativen Naturbetrachtung, die in der Stahlschen Phlogistontheorie noch einmal in charakteristischer Schärfe hervortritt. Aber als dann durch Lavoisiers Versuche der Umschwung eintritt, als auch alle chemischen Tatsachen unter den Gesichtpunkt von Maß, Zahl und Gewicht gestellt werden – da ist auch hier die Stunde für die Atomistik gekommen. Lavoisier erklärt, daß alle chemische Experimentierkunst auf dem Grundsatz beruhe, |
%%178	daß eine wahrhaftige Gleichheit oder Gleichung zwischen den anfänglich gegebenen Körpern und denen, die man aus ihnen durch die Analyse gewinnen kann, besteht.[73] So müssen auch alle Aussagen über die Atome, zu denen diese Experimentierkunst führt, zuletzt dazu dienen, dieser »Gleichheit« einen bestimmten Ausdruck zu geben und sie theoretisch verständlich zu machen. In Daltons Gesetz der multiplen Proportionen feiert die Atomistik ihren ersten Triumph; sie vermag unmittelbar die empirische Tatsache der konstanten Verbindungszahlen zu erklären. Die festen Äquivalentzahlen fanden ihre einfachste und sich unmittelbar aufdrängende Deutung in der Annahme, daß die Atome der verschiedenen Elemente ein bestimmtes Gewicht besitzen, das für die einzelnen Grundstoffe charakteristisch ist. Hierzu tritt | die Avogadrosche Regel, nach der die Anzahl der Moleküle in der Raumeinheit bei allen Gasen die gleiche ist. Auf dem einfachen Grundschema, das damit gewonnen ist, baut sich die gesamte folgende Entwicklung der Chemie auf.[74] Wie bei Huyghens, so wird hier das Atom streng substantiell gefaßt; aber alle Aussagen über sein »Was« gehen zuletzt in Aussagen über ein »Wieviel«, in Aussagen über bestimmte numerische Verhältnisse über.
%
%
%
%%For Perrin and Brownian motion, different experimental determinations converge on the same NA. Given the molar mass M, this also determines the mass of an individual molecule:
%%M m =
%%NA
%%So the connection to your previous point is direct: convergence fixes the numerical size of the molecular unit.
%%But convergence and discreteness remain distinct. Agreement on NA supports a theory in which matter consists of discrete molecules; the mere agreement of numerical results does not, by itself, establish discreteness.
%%On your Cassirerian reading, one could say: what we call the "molecule" receives its determinate scientific content through these convergent quantitative relations. Brownian motion supplies a further empirical route to that determination.
%
%
%
%%The difference is between substances supplying the determinations and methods supplying them. That does not by itself amount to two philosophically distinct kinds of convergence. Your broader point stands: Cassirer could understand Perrin’s results through the same pattern of numerical agreement across diverse empirical cases. One may interpret the convergence of independent experimental determinations as evidence for atoms. For Cassirer, however, what gives this interpretation determinate scientific content is precisely the system of numerical relations exhibited by that convergence.
%
%%Both Perrin’s argument and Cassirer’s examples concern numerical agreement across different empirical domains. Whether this agreement takes the form of independent determinations of the same constant or a precise relation between different constants, its evidential force lies in the coherence of otherwise independent results
%
%%Wie bei Huyghens, so wird hier das Atom streng substantiell gefaßt; aber alle Aussagen über sein »Was« gehen zuletzt in Aussagen über ein »Wieviel«, in Aussagen über bestimmte numerische Verhältnisse über. Alle Grundbegriffe und alle Theorien der wissenschaftlichen Chemie – der Valenzbegriff und der Typenbegriff wie der Begriff des Radikals und die Theorie der zusammengesetzten Radikale – weisen in diese Richtung. Die Begriffssprache der Atomistik erlaubt jetzt, in Gemeinschaft mit der allgemeinen Zeichensprache der Mathematik, jenen einfachen und großartigen Aufbau der chemischen Formelsprache, wie er sich im Fortgang von den einfachen Formeln zu den Strukturformeln vollzieht. Die Stereochemie, wie sie von van’t Hoff begründet wird, erlaubt auf Grund einer komplexeren Gestalt der Formeln eine neue Gruppe von Tatsachen, insbesondere die Erscheinungen der Isomerie, zum Ausdruck zu bringen; sie gibt einen neuen Aufschluß über die Lagerung der Atome im Raume. So wandelt und verfeinert sich das anschauliche Bild des Atoms; aber alle diese Wandlungen dienen nur dem einen Zweck, einen immer größeren Kreis funktioneller Beziehungen zu umspannen und darzustellen. 
%
%



\citetrackerfalse

%\printshorthands
\printbibliography

\end{document}